%% This is a LaTeX document.  Hey, Emacs, -*- latex -*- , get it?
\newif\ifbeamermode
\beamermodetrue
\ifbeamermode
\documentclass[mathserif,a4paper,aspectratio=169]{beamer}
\else
\documentclass[a4paper]{article}
\usepackage{beamerarticle}
\usepackage[a4paper,margin=2.5cm]{geometry}
\fi
\usepackage[shorthands=off,french]{babel}
\usepackage[utf8]{inputenc}
\usepackage[T1]{fontenc}
\usepackage{lmodern}
\DeclareUnicodeCharacter{00A0}{~}
\DeclareUnicodeCharacter{2026}{...}
\DeclareUnicodeCharacter{1E25}{\d{h}}
% Beamer theme:
\usetheme{Goettingen}
%\usecolortheme{albatross}
%\usecolortheme{lily}
%\setbeamercovered{transparent}
% A tribute to the worthy AMS:
\usepackage{amsmath}
\usepackage{amsfonts}
\usepackage{amssymb}
\usepackage{amsthm}
%
\usepackage{mathrsfs}
%
\usepackage{graphicx}
\usepackage{tikz}
\usetikzlibrary{arrows,automata,calc}
\usepackage{hyperref}
%
\newcommand{\itempoint}{\strut\hbox{\color{beamerstructure}\donotcoloroutermaths$\blacktriangleright$}\nobreak\hskip.5em plus.5em\relax}
\renewcommand{\thefootnote}{\textdagger}
\newcommand{\dbllangle}{\mathopen{\langle\!\langle}}
\newcommand{\dblrangle}{\mathclose{\rangle\!\rangle}}
%
\mode<article>{
    \newcommand{\donotcoloroutermaths}{\relax}
    \renewcommand{\frametitle}[1]{\subsubsection*{#1 {\footnotesize[transp. \insertframenumber]}}}
    \newcommand{\myitemmark}{\hbox{$\blacktriangleright$}}
    \renewcommand{\itempoint}{\strut\myitemmark\nobreak\hskip.5em plus.5em\relax}
    \setlength{\parindent}{0pt}
    \renewcommand{\thepage}{I–\arabic{page}}
}
%
%
%
\title{{\footnotesize Partie I:} Calculabilité}
\subtitle{CSC-3TC34-TP / INF110 (Logique et Fondements de l'Informatique)}
\author[David Madore]{David A. Madore\\
{\footnotesize Télécom Paris}\\
\texttt{david.madore@}(nom de l'école)\texttt{.fr}}
\date{2023–2026}
\mode<presentation>{%
\beamertemplatenavigationsymbolsempty
\usenavigationsymbolstemplate{\vbox{\hbox{\footnotesize(I) \hyperlinkslideprev{$\leftarrow$}\insertframenumber/\inserttotalframenumber\hyperlinkslidenext{$\rightarrow$}}}}
}
\setbeamercolor{myhighlight}{fg=black,bg=white!90!green}
\begin{document}
\mode<article>{\maketitle\renewcommand{\labelitemi}{\myitemmark}}
%
\setlength\abovedisplayskip{2pt plus 2pt minus 2pt}
\setlength\belowdisplayskip{2pt plus 2pt minus 2pt}
%
\begin{frame}
\titlepage
{\footnotesize\center{\url{http://perso.enst.fr/madore/inf110/transp-inf110.pdf}}\par}
{\tiny
\immediate\write18{sh ./vc > vcline.tex}
\begin{center}
Git: \input{vcline.tex}
\end{center}
\immediate\write18{echo ' (stale)' >> vcline.tex}
\par}
\end{frame}
%
\section*{Plan}
\begin{frame}
\frametitle{Plan}
\tableofcontents
\end{frame}
%
\section{Introduction}
\begin{frame}
\frametitle{Qu'est-ce que la calculabilité ?}

\itempoint À l'interface entre \textbf{logique mathématique} et
\textbf{informatique théorique}
\begin{itemize}
\item née de préoccupations venues de la logique (Hilbert, Gödel),
\item à l'origine des 1\textsuperscript{ers} concepts informatiques
  ($\lambda$-calcul, machine de Turing).
\end{itemize}

\bigskip

\itempoint But : étudier les limites de ce que \textbf{peut ou ne peut
  pas faire un algorithme}
\begin{itemize}
\item sans limite de ressources (temps, mémoire juste « finis »),
\item sans préoccupation d'efficacité ($\neq$ complexité, algorithmique),
\item y compris résultats négatifs (« \emph{aucun} algorithme ne peut… »),
\item voire relatifs (calculabilité relative),
\item admettant diverses généralisations (calculabilité supérieure).
\end{itemize}

\end{frame}
%
\begin{frame}
\frametitle{Quelques noms}

\itempoint Muḥammad ibn Mūsá al-\b{H}wārizmī (v.780–v.850) :
$\rightsquigarrow$« algorithme »

\itempoint Blaise Pascal (1623–1662) : machine à calculer
$\rightsquigarrow$automates

\itempoint Charles Babbage (1791–1871) : \textit{Analytical Engine} (Turing-complète !)

\itempoint Ada (née Byron) Countess of Lovelace (1815–1852) : programmation

\itempoint Richard Dedekind (1831–1916) : définitions primitives récursives

\itempoint David Hilbert (1862–1943) : \textit{Entscheidungsproblem}
(décider la vérité)

\itempoint Jacques Herbrand (1908–1931) : fonctions générales récursives

\itempoint Kurt Gödel (1906–1978) : incomplétude en logique

\itempoint Haskell Curry (1900–1982) : logique combinatoire, lien preuves-typage

\itempoint Alonzo Church (1903–1995) : $\lambda$-calcul

\itempoint Alan M. Turing (1912–1954) : machine de Turing, problème de l'arrêt

\itempoint Emil Post (1897–1954) : ensembles calculablement énumérables

\itempoint Stephen C. Kleene (1909–1994) : $\mu$-récursion, th. de récursion, forme normale

\end{frame}
%
\begin{frame}
\frametitle{Fonction calculable}

« Définition » : une fonction $f$ est \textbf{calculable}
quand il existe un algorithme qui
\begin{itemize}
\item prenant en entrée un $x$ du domaine de définition de $f$,
\item \textbf{termine en temps fini},
\item et renvoie la valeur $f(x)$.
\end{itemize}

\bigskip

Difficultés :
\begin{itemize}
\item Comment définir ce qu'est un algorithme ?
\item Quel type de valeurs (acceptées et renvoyées) ?
\item Et si l'algorithme ne termine pas ?
\item Distinction entre intention (l'algorithme) et extension (la fonction).
\end{itemize}

\end{frame}
%
\begin{frame}
\frametitle{Sans préoccupation d'efficacité}

\itempoint La calculabilité \alert{ne s'intéresse pas à l'efficacité}
des algorithmes qu'elle étudie, uniquement leur \textbf{terminaison en
  temps fini}.

\medskip

P.ex. : pour savoir si $n$ est premier, on peut tester si $i\times
j=n$ pour tout $i$ et $j$ allant de $2$ à $n-1$.  (Hyper inefficace ?
On s'en fiche.)

\bigskip

\itempoint La calculabilité \alert{n'a pas peur des grands entiers}.

\medskip

P.ex. : \textbf{fonction d'Ackermann} définie par :
\[
\begin{aligned}
A(m,n,0) &= m+n \\
A(m,1,k+1) &= m \\
A(m,n+1,k+1) &= A(m,\,A(m,n,k+1),\,k)
\end{aligned}
\]
définition algorithmique (par appels récursifs qui terminent), donc calculable.

\smallskip

Mais $A(2,6,3) = 2^{2^{2^{2^{2^2}}}} = 2^{2^{65\,536}}$ et $A(2,4,4) =
A(2,65\,536,3)$ est inimaginablement grand (et que dire de
$A(100,100,100)$ ?).

$\Rightarrow$ Ingérable sur un vrai ordinateur.

\end{frame}
%
\begin{frame}
\frametitle{Approches de la calculabilité}

\itempoint Approche informelle : \textbf{algorithme = calcul
  finitiste} mené par un humain ou une machine, selon des instructions
précises, en temps fini, sur des données finies

\medskip

\itempoint Approche pragmatique : tout ce qui peut être fait sur un
langage de programmation « Turing-complet » (Python, Java, C, Caml…)
idéalisé
\begin{itemize}
\item sans limites d'implémentation (p.ex., entiers arbitraires !),
\item sans source de hasard ou de non-déterminisme.
\end{itemize}

\medskip

\itempoint Approches formelles, p.ex. :
\begin{itemize}
\item fonctions générales récursives (Herbrand-Gödel-Kleene),
\item $\lambda$-calcul (Church) ($\leftrightarrow$ langages fonctionnels),
\item machine de Turing (Turing),
\item machines à registres (Post…).
\end{itemize}

\bigskip

\itempoint\textbf{« Thèse » de Church-Turing} : \alert{tout ceci
  donne la même chose}.

\end{frame}
%
\begin{frame}
\frametitle{Thèse de Church-Turing}

\itempoint\textbf{Théorème} (Post, Turing) : les fonctions (disons
$\mathbb{N} \dasharrow \mathbb{N}$) \textbf{(1)} générales récursives,
\textbf{(2)} représentables en $\lambda$-calcul, et
\textbf{(3)} calculables par machine de Turing, coïncident toutes.

\smallskip

$\Rightarrow$ On parle de \alert{calculabilité au sens de Church-Turing}.

\bigskip

\itempoint\textbf{Observation} : tous les langages de programmation
informatiques généraux usuels, idéalisés, calculent aussi exactement
ces fonctions ($\rightarrow$ « Turing-complets »).

\bigskip

\itempoint\textbf{Thèse philosophique} : la calculabilité de C-T
définit précisément la notion d'algorithme finitiste.

\bigskip

\itempoint\textbf{Conjecture physique} : la calculabilité de C-T
correspond aux calculs réalisables mécaniquement dans l'Univers (en
temps/énergie finis mais illimités).

{\footnotesize $\uparrow$ (même avec un ordinateur quantique)}

\bigskip

Pour toutes ces raisons, le sujet mérite d'être étudié !

\end{frame}
%
\begin{frame}
\frametitle{Trois grandes approches}

On va décrire trois approches des (mêmes !) fonctions calculables au
sens de Church-Turing, et esquisser leur équivalence :

\medskip

\itempoint Les \textbf{fonctions générales récursives} sont
mathématiq\textsuperscript{t} plus commodes :
\begin{itemize}
\item « tout est un entier » (fonctions $\mathbb{N}^k \dasharrow
  \mathbb{N}$),
\item définition inductive, numérotation associée.
\end{itemize}

\medskip

\itempoint Les \textbf{machines de Turing} représentent des
ordinateurs très simples :
\begin{itemize}
\item travaillent sur une « bande » illimitée a priori (mémoire),
\item aspect algorithmique évident, plus proche d'un « vrai » ordinateur,
\item approche la plus commode pour la complexité (pas considérée ici).
\end{itemize}

\medskip

\itempoint Le \textbf{$\lambda$-calcul} pur non typé est un système
symbolique :
\begin{itemize}
\item proche des langages de program\textsuperscript{tion} fonctionnels
  (Lisp, Haskell, OCaml…),
\item plus facile à « programmer » réellement, mais nombreuses
  subtilités.
\end{itemize}

\end{frame}
%
\begin{frame}
\frametitle{Données finies}

Un algorithme travaille sur des \textbf{données finies}.

\medskip

Qu'est-ce qu'une « donnée finie » ?  Tout objet représentable
informatiquement : booléen, entier, chaîne de caractères, structure,
liste/tableau de ces choses, ou même plus complexe (p.ex., graphe).

\medskip

$\rightarrow$ Comment y voir plus clair ?

\bigskip

Deux approches opposées :
\begin{itemize}
\item\textbf{typage} : distinguer toutes ces sortes de données,
\item\textbf{codage de Gödel} : tout représenter comme des entiers !
\end{itemize}

\bigskip

Le typage est plus élégant, plus satisfaisant, plus proche de
l'informatique réelle, on en reparlera.

\smallskip

Le codage de Gödel simplifie l'approche/définition de la calculabilité
(on étudie juste des fonctions $\mathbb{N} \dasharrow \mathbb{N}$).

\end{frame}
%
\begin{frame}
\label{codage-de-goedel}
\frametitle{Codage de Gödel (« tout est un entier »)}

\itempoint Représenter \textbf{n'importe quelle donnée finie par un
  entier}.

\bigskip

\itempoint Codage des couples : par exemple,
\[
\langle m,n\rangle := m + \frac{1}{2}(m+n)(m+n+1)
\]
définit une bijection calculable $\mathbb{N}^2 \to \mathbb{N}$
(calculable dans les deux sens).

\bigskip

\itempoint Codage des listes finies : par exemple,
\[
\dbllangle a_0,\ldots,a_{k-1}\dblrangle
:= \langle a_0, \langle a_1, \langle\cdots,\langle a_{k-1},0\rangle+1\cdots\rangle+1\rangle+1
\]
définit une bijection calculable $\{\text{suites finies dans $\mathbb{N}$}\} \to \mathbb{N}$ {\footnotesize (avec $\dbllangle\dblrangle := 0$)}.

%%% def encode_pair(m,n): return m+(m+n)*(m+n+1)/2
%%% def encode(t):
%%%     if isinstance(t, list):
%%%         v=0
%%%         for x in reversed(t):
%%%             m = encode(x)
%%%             v = encode_pair(m,v) + 1
%%%         return v
%%%     else: return t

\bigskip

\itempoint Il sera aussi utile de représenter même les
\alert{programmes} par des entiers.

\bigskip

\itempoint Les détails précis du codage sont \textbf{sans importance}.

\bigskip

\itempoint\textcolor{orange}{Ne pas utiliser dans la vraie vie} (hors calculabilité) !

\end{frame}
%
\begin{frame}
\frametitle{Fonctions partielles}

\itempoint Même si on s'intéresse à des algorithmes qui
\textbf{terminent}, la définition de la calculabilité \alert{doit
  forcément} passer aussi par ceux qui ne terminent pas.

{\footnotesize (Aucun langage Turing-complet ne peut exprimer
  uniquement des algorithmes qui terminent toujours, à cause de
  l'indécidabilité du problème de l'arrêt.)\par}

\bigskip

\itempoint Lorsque l'algorithme censé calculer $f(n)$ ne termine pas,
on dira que $f$ n'est pas définie en $n$, et on notera $f(n)\uparrow$.
Au contraire, s'il termine, on note $f(n)\downarrow$.

\bigskip

\itempoint Notation : $f\colon \mathbb{N} \dasharrow \mathbb{N}$ :
une fonction $D \to \mathbb{N}$ définie sur une partie $D \subseteq
\mathbb{N}$.

\itempoint Notation : $f(n) \downarrow$ signifie « $n \in D$ », et $f(n)
\uparrow$ signifie « $n \not\in D$ ».

\itempoint Notation : $f(n) \downarrow = g(m)$ signifie
« $f(n)\downarrow$ et $g(m)\downarrow$ et $f(n) = g(m)$ ».

\itempoint Convention : $f(n) = g(m)$ signifie « $f(n)\downarrow$ ssi
$g(m)\downarrow$, et $f(n) = g(m)$ si $f(n)\downarrow$ ».
{\footnotesize (Certains préfèrent écrire $f(n) \simeq g(m)$ pour ça.)}

\medskip

\itempoint Convention : si $g_i(\underline{x})\uparrow$ pour un $i$,
on convient que
$h(g_1(\underline{x}),\ldots,g_k(\underline{x}))\uparrow$.

\medskip

\itempoint Terminologie : une fonction $f\colon \mathbb{N} \to
\mathbb{N}$ définie sur $\mathbb{N}$ est dite \textbf{totale}.

{\footnotesize Une fonction totale est un \alert{cas particulier} de
  fonction partielle !\par}

\end{frame}
%
\begin{frame}
\frametitle{Terminologie à venir (avant-goût)}

\itempoint Une fonction partielle $f\colon \mathbb{N} \dasharrow
\mathbb{N}$ est dite \textbf{calculable} (partielle) lorsqu'il existe
un algorithme qui prend $n$ en entrée et :
\begin{itemize}
\item termine (en temps fini) et renvoie $f(n)$ lorsque $f(n)\downarrow$,
\item ne termine pas lorsque $f(n)\uparrow$.
\end{itemize}

\bigskip

\itempoint Une partie $A \subseteq \mathbb{N}$ est dite
\textbf{décidable} lorsque sa fonction indicatrice
$\mathbb{N}\to\mathbb{N}$
\[
\mathbf{1}_A\colon n \mapsto \left\{
\begin{array}{ll}
1&\text{~si~}n\in A\\
0&\text{~si~}n\not\in A\\
\end{array}
\right.
\]
est calculable (répondre « oui » ou « non » selon que $n\in A$ ou $n\not\in A$).

\bigskip

\itempoint Une partie $A \subseteq \mathbb{N}$ est dite
\textbf{semi-décidable} lorsque sa fonction partielle « semi-indicatrice »
$\mathbb{N}\dasharrow\mathbb{N}$ (d'ensemble de définition $A$)
\[
n \mapsto \left\{
\begin{array}{ll}
1&\text{~si~}n\in A\\
\uparrow&\text{~si~}n\not\in A\\
\end{array}
\right.
\]
est calculable (répondre « oui » ou « ... » selon que $n\in A$ ou $n\not\in A$).

\end{frame}
%
\begin{frame}
\frametitle{Point terminologique : « récursif »}

Le mot « récursif » et ses cognats (« récursion », « récursivité ») a
plusieurs sens \alert{apparentés mais non identiques} :
\begin{itemize}
\item « récursif » = « défini par récurrence » (Dedekind 1888)
  $\rightarrow$ fonctions primitives récursives, générales
  récursives (cf. après) ;
\item « récursif » = « calculable » (par glissement à cause de la
  définition de la calculabilité par les fonctions générales
  récursives) ;
\item « récursif » = « faisant appel à lui-même dans sa définition »
  (appels récursifs, récursivité en informatique).
\end{itemize}

\bigskip

On va définir les fonctions « \textbf{primitives récursives} »
(1\textsuperscript{er} sens) et « \textbf{(générales) récursives} »
(1\textsuperscript{er} et aussi 2\textsuperscript{e} sens) ci-après.

\medskip

Pour le 3\textsuperscript{e} sens, on dira « appels récursifs ».

\end{frame}
%
\section{Fonctions primitives récursives}
\begin{frame}
\frametitle{Fonctions primitives récursives : aperçu}

\itempoint Avant de définir les fonctions générales récursives
($\cong$ calculables), on va commencer par les \textbf{primitives
  récursives}, plus restreintes.

{\footnotesize« primitive\alert{ment} récursives » ?\par}

\bigskip

\itempoint Historiquement antérieures à la calculabilité de
Church-Turing.

\bigskip

\itempoint Pédagogiquement utile comme « échauffement ».

\bigskip

\itempoint À cheval entre calculabilité (\textbf{PR} est une petite
classe de calculabilité) et complexité (c'est une grosse classe de
complexité).

\bigskip

\itempoint Correspond à des programmes à \textbf{boucles bornées a
  priori}.

\bigskip

\itempoint Énormément d'algorithmes usuels sont p.r.

\bigskip

\itempoint Mais pas tous : p.ex. la fonction d'Ackermann n'est pas p.r.

\end{frame}
%
\begin{frame}
\label{primitive-recursive-definition}
\frametitle{Fonctions primitives récursives : définition}

\itempoint $\textbf{PR}$ est la plus petite classe de fonctions
$\mathbb{N}^k \dasharrow \mathbb{N}$ (en fait $\mathbb{N}^k \to
\mathbb{N}$), pour $k$ variable qui :
\begin{itemize}
\item contient les projections $\underline{x} := (x_1,\ldots,x_k)
  \mapsto x_i$ ;
\item contient les constantes $\underline{x} \mapsto c$ ;
\item contient la fonction successeur $x \mapsto x+1$ ;
\item est stable par composition : si $g_1,\ldots,g_\ell\colon
  \mathbb{N}^k \dasharrow \mathbb{N}$ et $h\colon \mathbb{N}^\ell
  \dasharrow \mathbb{N}$ sont p.r. alors $\underline{x} \mapsto
  h(g_1(\underline{x}),\ldots, g_\ell(\underline{x}))$ est p.r. ;
\item est stable par récursion primitive : si $g\colon \mathbb{N}^k
  \dasharrow \mathbb{N}$ et $h\colon \mathbb{N}^{k+2} \dasharrow
  \mathbb{N}$ sont p.r., alors $f\colon \mathbb{N}^{k+1} \dasharrow
  \mathbb{N}$ est p.r., où :
\[
\begin{aligned}
f(\underline{x},0) &= g(\underline{x})\\
f(\underline{x},z+1) &= h(\underline{x},f(\underline{x},z),z)
\end{aligned}
\]
\end{itemize}

\medskip

{\footnotesize Les fonctions p.r. sont automatiq\textsuperscript{t}
  totales, mais il est commode de garder la définition avec
  $\dasharrow$.\par}

\end{frame}
%
\begin{frame}
\frametitle{Fonctions primitives récursives : exemples}

\itempoint $f\colon (x,z) \mapsto x+z$ est p.r. :
\[
\begin{aligned}
f(x,0) &= x\\
f(x,z+1) &= f(x,z)+1
\end{aligned}
\]
{\footnotesize où $x \mapsto x$ et $(x,y,z) \mapsto y+1$ sont p.r.\par}

\medskip

\itempoint $f\colon (x,z) \mapsto x\cdot z$ est p.r. :
\[
\begin{aligned}
f(x,0) &= 0\\
f(x,z+1) &= f(x,z)+x
\end{aligned}
\]

\medskip

\itempoint $f\colon (x,z) \mapsto x^z$ est p.r.

\bigskip

\itempoint $f\colon (x,y,0) \mapsto x, \; (x,y,z) \mapsto y\text{~si~}z\geq 1$ est p.r. :
\[
\begin{aligned}
f(x,y,0) &= x\\
f(x,y,z+1) &= y
\end{aligned}
\]

\medskip

\itempoint $u \mapsto \max(u-1,0)$ est p.r. (exercice !), comme $(u,v)
\mapsto \max(u-v,0)$ ou $(u,v) \mapsto u\% v$ ou $(u,v) \mapsto
\lfloor u/v\rfloor$.

\end{frame}
%
\begin{frame}
\frametitle{Fonctions primitives récursives : programmation}

Les fonctions p.r. sont celles définies par un \textbf{langage de
  programmation à boucles bornées}, c'est-à-dire que :
\begin{itemize}
\item les variables sont des entiers naturels (illimités !),
\item les manipulations de base sont permises (constantes,
  affectations, test d'égalité, conditionnelles),
\item les opérations arithmétiques basiques sont disponibles,
\item on peut faire des appels de fonctions \alert{sans appels récursifs},
\item on ne peut faire que des boucles \alert{de nombre borné
  \textit{a priori}} d'itérations.
\end{itemize}

\medskip

Les programmes dans un tel langage \textbf{terminent forcément par
  construction}.

\bigskip

\textbf{N.B.} $(m,n) \mapsto \langle m,n\rangle := m +
\frac{1}{2}(m+n)(m+n+1)$ et $\langle m,n\rangle \mapsto m$ et $\langle
m,n\rangle \mapsto n$ sont p.r.

\end{frame}
%
\begin{frame}
\frametitle{Fonctions primitives récursives : lien avec la complexité}

En anticipant sur la notion de machine de Turing :

\medskip

\itempoint La fonction $(M,C) \mapsto C'$ qui à une machine de Turing
$M$ et une configuration (= ruban+état) $C$ de $M$ associe la
configuration suivante \textbf{est p.r.}

\medskip

\itempoint Conséquence : la fonction $(n,M,C) \mapsto C^{(n)}$ qui à
$n\in\mathbb{N}$ et une machine de Turing $M$ et une configuration $C$
de $M$ associe la configuration atteinte après $n$ étapes d'exécution,
\textbf{est p.r.}

{\footnotesize (Par récursion primitive sur le point précédent.)}

\medskip

\itempoint Conséquence : une fonction calculable en complexité
p.r. par une machine de Turing est elle-même p.r.

\smallskip

{\footnotesize (Calculer une borne p.r. sur le nombre d'étapes, puis
  appliquer le point précédent.)}

\medskip

\itempoint Réciproquement : une p.r. est calculable en complexité p.r.

\medskip

\itempoint Moralité : p.r. $\Leftrightarrow$ de complexité p.r.

\smallskip

{\footnotesize Notamment $\textbf{EXPTIME} \subseteq \textbf{PR}$.\par}

\end{frame}
%
\begin{frame}
\frametitle{Fonctions primitives récursives : limitations}

{\footnotesize La classe $\textbf{PR}$ est « à cheval » entre la
  calculabilité et la complexité.\par}

\bigskip

Rappel : la \textbf{fonction d'Ackermann} (pour $m=2$) définie par :
\[
\begin{aligned}
A(2,n,0) &= 2+n \\
A(2,1,k+1) &= 2 \\
A(2,n+1,k+1) &= A(2,\,A(2,n,k+1),\,k)
\end{aligned}
\]
devrait être calculable.  Mais cette définition \alert{n'est pas une
  récursion primitive} (pourquoi ?).

\bigskip

\itempoint On peut montrer que : si $f \colon \mathbb{N}^k \to
\mathbb{N}$ est p.r., il existe $r$ ($=r(f)$) tel que
\[
f(x_1,\ldots,x_k) \leq A(2,\, (x_1+\cdots+x_k+3),\, r)
\]

\medskip

\itempoint Notamment, $r \mapsto A(2, r, r)$ \textbf{n'est pas p.r.}

\medskip

Pourtant, \alert{elle est bien définie par un algorithme} clair (et
terminant clairement).

\end{frame}
%
\begin{frame}
\frametitle{Fonctions primitives récursives : numérotation (idée)}

\itempoint On veut \alert{coder} les fonctions p.r. {\footnotesize (et
  plus tard : gén\textsuperscript{ales} récursives)} \alert{par des
  entiers}.

\bigskip

\itempoint Pour (certains) entiers $e \in \mathbb{N}$, on va définir
$\psi_e^{(k)}\colon \mathbb{N}^k \to \mathbb{N}$ primitive récursive,
la fonction p.r. \alert{codée} par $e$ ou ayant $e$ comme
\textbf{code} (source) / « programme ».

\bigskip

\itempoint Toute fonction p.r. $f\colon \mathbb{N}^k \to \mathbb{N}$
sera un $\psi_e^{(k)}$ pour un certain $e$.

\smallskip

\itempoint Ce $e$ décrit la manière dont $f$ est construite selon la
définition de $\mathbf{PR}$
(cf. transp. \ref{primitive-recursive-definition}).

\smallskip

\itempoint Il faut l'imaginer comme le \alert{code source} de $f$ (au
sens informatique).

\smallskip

\itempoint Il n'est \alert{pas du tout unique} : $f = \psi_{e_1}^{(k)}
= \psi_{e_2}^{(k)} = \cdots$

{\footnotesize ($e$ = « intention » / $f$ = « extension »)}

\bigskip

{\footnotesize

\itempoint On va ensuite se demander si $(e,\underline{x}) \mapsto
\psi_e^{(k)}(\underline{x})$ est \alert{elle-même p.r.} (divulgâchis :
\alert{non}).

\par}

\end{frame}
%
\begin{frame}
\frametitle{Fonctions primitives récursives : numérotation (définition)}

On définit $\psi_e^{(k)}\colon \mathbb{N}^k \dasharrow \mathbb{N}$ par
induction suivant la déf\textsuperscript{n} de $\mathbf{PR}$
(cf. transp. \ref{primitive-recursive-definition}) :
\begin{itemize}
\item si $e = \dbllangle 0, k, i\dblrangle$ alors
  $\psi_e^{(k)}(x_1\ldots,x_k) = x_i$ (projections) ;
\item si $e = \dbllangle 1, k, c\dblrangle$ alors
  $\psi_e^{(k)}(x_1\ldots,x_k) = c$ (constantes) ;
\item si $e = \dbllangle 2\dblrangle$ alors
  $\psi_e^{(1)}(x) = x+1$ (successeur) ;
\item si $e = \dbllangle 3, k, d, c_1,\ldots,c_\ell\dblrangle$ et $g_i
  := \psi_{c_i}^{(k)}$ et $h := \psi_d^{(\ell)}$, alors
  $\psi_e^{(k)} \colon \underline{x} \mapsto
  h(g_1(\underline{x}),\ldots, g_\ell(\underline{x}))$ (composition) ;
\item si $e = \dbllangle 4, k, d, c\dblrangle$ et $g :=
  \psi_c^{(k)}$ et $h := \psi_d^{(k+2)}$, alors (récursion primitive)
\[
\begin{aligned}
\psi_e^{(k+1)}(\underline{x},0) &= g(\underline{x})\\
\psi_e^{(k+1)}(\underline{x},z+1) &= h(\underline{x}, \psi_e^{(k+1)}(\underline{x},z),z)
\end{aligned}
\]
\end{itemize}
(Autres cas non définis, i.e., donnent $\uparrow$.)

\bigskip

\itempoint Alors $f\colon \mathbb{N}^k \dasharrow \mathbb{N}$ est
p.r. \alert{ssi} $\exists e \in\mathbb{N}.\,(f = \psi_e^{(k)})$
par définition.

{\tiny P.ex., $e = \dbllangle 4,1,\dbllangle 3,3,\dbllangle
  2\dblrangle,\dbllangle 0,3,2\dblrangle\dblrangle,\dbllangle
  0,1,1\dblrangle\dblrangle =
  1\,459\,411\,784\,487\,\ldots\,780\,615\,609\,825 \approx
  1.459\times 10^{357}$ définit $\psi^{(2)}_e(x,z) = x+z$ avec les
  conventions de codage du transp. \ref{codage-de-goedel}.\par}

\end{frame}
%
\begin{frame}
\frametitle{Manipulation de programmes (version p.r.)}

\itempoint Penser à $e$ dans $\psi_e^{(k)}$ comme un programme écrit
en « langage p.r. ».

\medskip

\itempoint La fonction $\psi_e^{(k)}\colon \mathbb{N}^k \dasharrow
\mathbb{N}$ « interprète » le programme $e$.

\medskip

\itempoint Une fonction p.r. donnée a \alert{beaucoup d'indices} :
$\psi_{e_1}^{(k)} = \psi_{e_2}^{(k)} = \cdots$ (programmes équivalents).

\medskip

\centerline{*}

\bigskip

La numérotation (transp. précédent) rend p.r. beaucoup de
manipulations usuelles de programmes (composition, récursion, etc.).
Notamment :

\medskip

\itempoint\textbf{Théorème s-m-n} (Kleene) : il existe $s_{m,n} \colon
\mathbb{N}^{m+1} \to \mathbb{N}$ p.r. telle que
\[
(\forall e,\underline{x},\underline{y})\quad
\psi^{(n)}_{s_{m,n}(e,x_1,\ldots,x_m)}(y_1,\ldots,y_n) =
\psi^{(m+n)}_e(x_1,\ldots,x_m,\,y_1,\ldots,y_n)
\]

{\footnotesize\underline{Preuve :} $s_{m,n}(e,\underline{x}) =
  \dbllangle 3, n, e, \dbllangle 1, n, x_1\dblrangle, \ldots,
  \dbllangle 1, n, x_m\dblrangle, \; \dbllangle 0, n, 1\dblrangle,
  \ldots, \dbllangle 0, n, n\dblrangle \dblrangle$ avec nos
  conventions (composition de fonctions constantes et de
  projections).\qed\par}

\medskip

\emph{En clair :} $s_{m,n}$ prend un programme $e$ qui prend $m+n$
arguments en entrée et « fixe » la valeur des $m$ premiers arguments à
$x_1,\ldots,x_m$, les $n$ arguments suivants ($y_1,\ldots,y_n$) étant
gardés variables.

\end{frame}
%
\begin{frame}
\frametitle{Digression : l'astuce de Quine (intuition)}

{\footnotesize Le nom de Willard Van Orman Quine (1908–2000) a été
  associé à cette astuce par Douglas Hofstadter.  En fait, l'astuce
  est plutôt due à Cantor, Gödel, Turing ou Kleene.\par}

\smallskip

\textcolor{teal}{Les mots suivants suivis des mêmes mots entre
  guillemets forment une phrase intéressante : « les mots suivants
  suivis des mêmes mots entre guillemets forment une phrase
  intéressante ».}

\bigskip

Pseudocode :

\smallskip

{\footnotesize\texttt{%
str="somefunc(code) \{ /*...*/ \}\textbackslash nsomefunc(\textbackslash"str=\textbackslash"+quote(str)+str);\textbackslash n";\newline
somefunc(code) \{ /*...*/ \}\newline
somefunc("str="+quote(str)+str);
}\par}

\smallskip

$\Rightarrow$ La fonction \texttt{somefunc} (arbitraire) est appelée
avec le code source du programme \alert{tout entier}.

\medskip

{\footnotesize\textbf{Exercice :} utiliser cette astuce pour écrire un
  programme écrivant son propre code source.\par}

\bigskip

\textcolor{blue}{\textbf{Moralité :}} \alert{on peut toujours donner aux programmes
  accès à leur code source}, même si ce n'est pas prévu par le
langage.

\end{frame}
%
\begin{frame}
\label{kleene-recursion-theorem-p-r-version}
\frametitle{Le théorème de récursion de Kleene (version p.r.)}

Version formelle de l'astuce de Quine

\smallskip

\itempoint\textbf{Théorème} (Kleene) : si $h \colon \mathbb{N}^{k+1}
\dasharrow \mathbb{N}$ est p.r., il existe $e$ tel que
\[
(\forall\underline{x})\quad \psi^{(k)}_e(\underline{x}) = h(e,\underline{x})
\]
Plus précisément, il existe $b \colon \mathbb{N}^2 \to \mathbb{N}$
p.r. telle que $e := b(k,d)$ vérifie
\[
(\forall d)\,(\forall\underline{x})\quad \psi^{(k)}_e(\underline{x}) = \psi^{(k+1)}_d(e,\underline{x})
\]

\bigskip

\underline{Preuve :} soit $s := s_{1,k}$ donné par le théorème s-m-n.
La fonction $(t,\underline{x}) \mapsto h(s(t,t),\underline{x})$ est
p.r., disons $= \psi_c^{(k+1)}(t,\underline{x})$.  Alors
\[
\psi_{s(c,c)}^{(k)}(\underline{x})
= \psi_{c}^{(k+1)}(c, \underline{x})
= h(s(c,c),\underline{x})
\]
donc $e := s(c,c)$ convient.  Les fonctions $d \mapsto c$ et $c \mapsto e$
sont p.r.\qed

\bigskip

\textcolor{blue}{\textbf{Moralité :}} \alert{on peut donner aux
  programmes accès à leur propre numéro} (= « code source », ici $e$),
cela ne change rien.

\end{frame}
%
\begin{frame}
\label{primitive-recursive-no-universality}
\frametitle{Fonctions primitives récursives : absence d'universalité}

\itempoint\textbf{Théorème :} il n'existe pas de fonction
p.r. $u\colon \mathbb{N}^2 \to \mathbb{N}$ telle que $u(e,x) =
\psi^{(1)}_e(x)$ si $\psi^{(1)}_e(x)\downarrow$.

\bigskip

\underline{Preuve :} par l'absurde : si un tel $u$ existe, alors
$(e,x) \mapsto u(e,x)+1$ est p.r.  Par le théorème de récursion de
Kleene, il existe $e$ tel que $\psi^{(1)}_e(x) = u(e,x) + 1$, ce qui
contredit $u(e,x) = \psi^{(1)}_e(x)$.\qed

\medskip

\centerline{*}

\medskip

\textcolor{blue}{\textbf{Moralité :}} \alert{un interpréteur du
  langage p.r. ne peut pas être p.r.} (preuve : on peut interpréter
l'interpréteur s'interprétant lui-même, en ajoutant $1$ au résultat
ceci donne un paradoxe ; c'est un argument diagonal de Cantor).

\bigskip

\itempoint Cet argument dépend du théorème s-m-n et du fait que les
fonctions p.r. sont \alert{totales}.  Pour définir une théorie
satisfaisante de la calculabilité, on va sacrifier la totalité pour
sauver le théorème s-m-n.

{\footnotesize Cette même preuve deviendra alors la preuve de
  l'indécidabilité du problème de l'arrêt.\par}

\end{frame}
%
\begin{frame}
\frametitle{Fonctions primitives récursives : absence d'universalité (variante)}

Rappel : la \textbf{fonction d'Ackermann} est définie par :
\[
\begin{aligned}
A(m,n,0) &= m+n \\
A(m,1,k+1) &= m \\
A(m,n+1,k+1) &= A(m,\,A(m,n,k+1),\,k)
\end{aligned}
\]

\bigskip

\itempoint Pour un $k$ \alert{fixé}, la fonction $(m,n) \mapsto
A(m,n,k)$ est p.r. (par récurrence sur $k$, récursion primitive sur
$A(m,n,k-1)$).

\bigskip

\itempoint Il existe même $k \mapsto a(k)$ p.r. telle que
$\psi^{(2)}_{a(k)}(m,n) = A(m,n,k)$.

\smallskip

I.e., on peut calculer de façon p.r. en $k$ le \alert{code} d'un
programme p.r. qui calcule $(m,n) \mapsto A(m,n,k)$.

\bigskip

\itempoint Si (une extension de) $(e,n) \mapsto \psi^{(1)}_e(n)$ était
p.r., on pourrait calculer $(n,k) \mapsto
\psi^{(1)}_{s_{1,1}(a(k),2)}(n) = \psi^{(2)}_{a(k)}(2,n) = A(2,n,k)$,
or elle n'est pas p.r.

\end{frame}
%
\section{Fonctions générales récursives}
\begin{frame}
\frametitle{Fonctions générales récursives : aperçu}

\itempoint On a vu que les fonctions p.r. sont \alert{limitées} et ne
couvrent pas la notion générale d'algorithme :
\begin{itemize}
\item les algorithmes p.r. terminent toujours car
\item le langage ne permet pas de boucles non bornées ;
\item concrètement, il n'implémente pas la fonction d'Ackermann ;
\item il ne peut pas s'interpréter lui-même.
\end{itemize}

\bigskip

\itempoint On veut modifier la définition des fonctions p.r. pour
lever ces limitations.  On va \alert{autoriser les boucles infinies}.

$\rightarrow$ Fonctions \textbf{générales récursives} ou simplement
« \textbf{récursives} ».

Ce seront aussi nos fonctions \textbf{calculables} !

\bigskip

\itempoint En ce faisant, on obtient forcément des cas de
non-terminaisons, donc on doit passer par des \alert{fonctions
  partielles}.

\bigskip

{\footnotesize\textbf{N.B.} Terminologie fluctuante : fonctions
  « générales récursives » ? juste « récursives » ? « récursives
  partielles » ?  « calculables » ? « calculables partielles » ?\par}

\end{frame}
%
\begin{frame}
\frametitle{L'opérateur $\mu$ de Kleene}

\textbf{Définition :} si $g\colon \mathbb{N}^{k+1} \dasharrow
\mathbb{N}$ et $\underline{x} \in \mathbb{N}^k$, alors $\mu
g(\underline{x})$ est le plus petit $z$ tel que $g(z,\underline{x}) =
0$ et $g(i,\underline{x})\downarrow$ pour $0\leq i<z$, s'il existe.

\bigskip

Autrement dit, $\mu g(\underline{x}) = z$ signifie
\begin{itemize}
\item $g(z,\underline{x}) = 0$,
\item $g(i,\underline{x})>0$ (sous-entendant
  $g(i,\underline{x})\downarrow$) pour tout $0\leq i<z$.
\end{itemize}

\bigskip

Concrètement, penser à $\mu g$ comme la fonction

\texttt{i=0;  while (true) \{ if (g(i,x)==0) \{ return i; \}  i++; \}}

\bigskip

\itempoint Ceci permet toute sorte de \alert{recherche non bornée}.

\bigskip

\itempoint L'opérateur $\mu$ associe à une fonction $g\colon
\mathbb{N}^{k+1} \dasharrow \mathbb{N}$ une fonction $\mu g\colon
\mathbb{N}^k \dasharrow \mathbb{N}$.

\end{frame}
%
\begin{frame}
\label{general-recursive-definition}
\frametitle{Fonctions générales récursives : définition}

\itempoint $\textbf{R}$ est la plus petite classe de fonctions
$\mathbb{N}^k \dasharrow \mathbb{N}$, pour $k$ variable qui :
\begin{itemize}
\item contient les projections $\underline{x} := (x_1,\ldots,x_k)
  \mapsto x_i$ ;
\item contient les constantes $\underline{x} \mapsto c$ ;
\item contient la fonction successeur $x \mapsto x+1$ ;
\item est stable par composition : si $g_1,\ldots,g_\ell\colon
  \mathbb{N}^k \dasharrow \mathbb{N}$ et $h\colon \mathbb{N}^\ell
  \dasharrow \mathbb{N}$ sont récursives alors $\underline{x} \mapsto
  h(g_1(\underline{x}),\ldots, g_\ell(\underline{x}))$ est récursive ;
\item est stable par récursion primitive : si $g\colon \mathbb{N}^k
  \dasharrow \mathbb{N}$ et $h\colon \mathbb{N}^{k+2} \dasharrow
  \mathbb{N}$ sont récursives, alors $f\colon \mathbb{N}^{k+1}
  \dasharrow \mathbb{N}$ est récursive, où :
\[
\begin{aligned}
f(\underline{x},0) &= g(\underline{x})\\
f(\underline{x},z+1) &= h(\underline{x},f(\underline{x},z),z)
\end{aligned}
\]
\item est stable par l'opérateur $\mu$ : si $g\colon \mathbb{N}^{k+1}
  \dasharrow \mathbb{N}$ est récursive, alors $\mu g\colon
  \mathbb{N}^k \dasharrow \mathbb{N}$ est récursive
  \textcolor{teal}{($\leftarrow$ nouveau !)}.
\end{itemize}

\medskip

Cette fois le langage \alert{permet les boucles non bornées} !

\end{frame}
%
\begin{frame}
\frametitle{Fonctions générales récursives : numérotation (idée)}

\itempoint On veut \alert{coder} les fonctions générales récursives
\alert{par des entiers}.

\smallskip

{\footnotesize Exactement comme on l'a fait pour les fonctions p.r.,
  on change juste la notation de $\psi$ en $\varphi$.\par}

\bigskip

\itempoint Pour (certains) entiers $e \in \mathbb{N}$, on va définir
$\varphi_e^{(k)}\colon \mathbb{N}^k \dasharrow \mathbb{N}$ générale
récursive, celle \alert{codée} par $e$ ou ayant $e$ comme
\textbf{code} (source) / « programme ».

\bigskip

\itempoint Toute fonction récursive (partielle !) $f\colon
\mathbb{N}^k \dasharrow \mathbb{N}$ sera un $\varphi_e^{(k)}$ pour un
certain $e$.

\smallskip

\itempoint Ce $e$ décrit la manière dont $f$ est construite selon la
définition de $\mathbf{R}$
(cf. transp. \ref{general-recursive-definition}).

\smallskip

\itempoint Il faut l'imaginer comme le \alert{code source} de $f$ (au
sens informatique).

\smallskip

\itempoint Il n'est \alert{pas du tout unique} : $f = \varphi_{e_1}^{(k)}
= \varphi_{e_2}^{(k)} = \cdots$

{\footnotesize ($e$ = « intention » / $f$ = « extension »)}

\bigskip

{\footnotesize

\itempoint On va ensuite se demander si $(e,\underline{x}) \mapsto
\varphi_e^{(k)}(\underline{x})$ est \alert{elle-même récursive}
(divulgâchis : \alert{oui}, \alert{contrairement} au cas p.r.).

\par}

\end{frame}
%
\begin{frame}
\frametitle{Fonctions générales récursives : numérotation (définition)}

On définit $\varphi_e^{(k)}\colon \mathbb{N}^k \dasharrow \mathbb{N}$ par
induction suivant la déf\textsuperscript{n} de $\mathbf{R}$
(cf. transp. \ref{general-recursive-definition}) :
\begin{itemize}
\item si $e = \dbllangle 0, k, i\dblrangle$ alors
  $\varphi_e^{(k)}(x_1\ldots,x_k) = x_i$ (projections) ;
\item si $e = \dbllangle 1, k, c\dblrangle$ alors
  $\varphi_e^{(k)}(x_1\ldots,x_k) = c$ (constantes) ;
\item si $e = \dbllangle 2\dblrangle$ alors
  $\varphi_e^{(1)}(x) = x+1$ (successeur) ;
\item si $e = \dbllangle 3, k, d, c_1,\ldots,c_\ell\dblrangle$ et $g_i
  := \varphi_{c_i}^{(k)}$ et $h := \varphi_d^{(\ell)}$, alors
  $\varphi_e^{(k)} \colon \underline{x} \mapsto
  h(g_1(\underline{x}),\ldots, g_\ell(\underline{x}))$ (composition) ;
\item si $e = \dbllangle 4, k, d, c\dblrangle$ et $g :=
  \varphi_c^{(k)}$ et $h := \varphi_d^{(k+2)}$, alors (récursion primitive)
\[
\begin{aligned}
\varphi_e^{(k+1)}(\underline{x},0) &= g(\underline{x})\\
\varphi_e^{(k+1)}(\underline{x},z+1) &= h(\underline{x},f(\underline{x},z),z)
\end{aligned}
\]
\item si $e = \dbllangle 5, k, c\dblrangle$ et $g :=
  \varphi_c^{(k+1)}$, alors $\varphi_e^{(k)} = \mu g$.
\end{itemize}
(Autres cas non définis, i.e., donnent $\uparrow$.)

\bigskip

\itempoint Alors $f\colon \mathbb{N}^k \dasharrow \mathbb{N}$ est
récursive \alert{ssi} $\exists e \in\mathbb{N}.\,(f = \varphi_e^{(k)})$
par définition.

\end{frame}
%
\begin{frame}
\frametitle{Fonctions générales récursives : universalité}

Cette fois, \alert{il existe bien} une fonction récursive universelle,
c'est-à-dire :

\smallskip

\itempoint $(e,\underline{x}) \mapsto \varphi_e^{(k)}(\underline{x})$
est récursive : $\exists u_k \in
\mathbb{N}.\,(\varphi_{u_k}^{(k+1)}(e, \underline{x}) =
\varphi_e^{(k)}(\underline{x}))$.

{\footnotesize Variante : $\exists u.\,(\varphi_{u}^{(1)}(\dbllangle
  e, \underline{x}\dblrangle) = \varphi_e^{(k)}(\underline{x}))$ en
  codant programme $e$ et arguments $\underline{x}$ en un entier.\par}

\bigskip

Qu'est-ce que ça nous dit ?

\medskip

\itempoint\textcolor{blue}{Mathématiquement}, que $\varphi^{(k)}$ est
elle-même récursive en tous ses arguments, y compris l'indice de
numérotation.

\medskip

\itempoint\textcolor{blue}{Informatiquement}, qu'il existe un
\alert{interpréteur} $u$ du langage général récursif dans le langage
général récursif.

\medskip

\itempoint\textcolor{blue}{Philosophiquement}, que la notion
d'« exécuter un algorithme » est \alert{elle-même algorithmique}.

\bigskip

\textcolor{brown}{Mais comment le prouver ?}  On va esquisser une
méthode par \textbf{arbres de calcul}.

\end{frame}
%
\begin{frame}
\frametitle{Arbres de calcul pour les fonctions récursives}

Un \textbf{arbre de calcul} de $\varphi_e^{(k)}(\underline{x})$ de
résultat $y$ est un arbre (fini, enraciné, ordonné, étiqueté par des
entiers naturels) qui « atteste » que $\varphi_e^{(k)}(\underline{x})
= y$ en détaillant les étapes du calcul.

\medskip

\begin{itemize}
\item La racine est étiquetée « $\varphi_e^{(k)}(\underline{x}) = y$ »
  codé disons « $\dbllangle e, \dbllangle \underline{x}\dblrangle,
  y\dblrangle$ ».
\item Le sous-arbre porté par chaque fils de la racine est lui-même un
  arbre de calcul pour une sous-expression utilisée dans le calcul de
  $f(\underline{x}) = y$.
\item Pour les cas projection, constante, successeur, il n'y a pas de
  fils.
\item Pour la composition
  $h(g_1(\underline{x}),\ldots,g_\ell(\underline{x}))$, les fils
  portent des arbres de calcul attestant
  $g_1(\underline{x})=v_1,\ldots,g_\ell(\underline{x})=v_\ell$ et
  $h(\underline{v})=y$ (donc $\ell+1$ fils).
\item Pour la récursion primitive $f(\underline{x},z)$, on a soit un
  seul fils attestant $g(\underline{x})=y$ lorsque $z=0$ soit deux
  fils attestant $f(\underline{x},z')=v$ et $h(\underline{x},v,z')=y$
  lorsque $z=z'+1$ (donc $1$ ou $2$ fils).
\item Pour $f = \mu g$, on a des fils attestant
  $g(0,\underline{x})=v_0,\ldots,g(y,\underline{x})=v_y$ où $v_i>0$ si
  $0\leq i<y$ et $v_y=0$ (donc $y+1$ fils).
\end{itemize}

\medskip

On a $\varphi_e^{(k)}(\underline{x}) = y$ \alert{ssi} il en existe un
arbre de calcul qui l'atteste.

\end{frame}
%
\begin{frame}
\frametitle{Arbres de calcul : formalisation}

Formellement :

{\footnotesize
\begin{itemize}
\item la racine est étiquetée $\dbllangle e, \dbllangle
  \underline{x}\dblrangle, y\dblrangle$,
\item soit $e$ est $\dbllangle 0, k, i\dblrangle$ ou $\dbllangle 1, k,
  c\dblrangle$ ou $\dbllangle 2\dblrangle$, elle n'a pas de
  fils, et $y$ vaut resp\textsuperscript{t} $x_i$, $c$ ou $x+1$,
\item soit $e = \dbllangle 3, k, d, c_1,\ldots,c_\ell\dblrangle$ et
  les $\ell+1$ fils portent des arbres de calculs attestant
  $\varphi_{c_1}^{(k)}(\underline{x})=v_1$, ...,
  $\varphi_{c_\ell}^{(k)}(\underline{x})=v_\ell$ et
  $\varphi_d^{(\ell)}(\underline{v})=y$.
\item soit $e = \dbllangle 4, k', d, c\dblrangle$ où $k'=k-1$ et
\begin{itemize}\footnotesize
\item soit $x_k = 0$ et l'unique fils porte arbre de calcul attestant
  $\varphi_{c}^{(k')}(x_1,\ldots,x_{k'})=y$,
\item soit $x_k = z+1$ et les deux fils portent arbres de calcul
  attestant $\varphi_{e}^{(k'+1)}(x_1,\ldots,x_{k'},z)=v$ et
  $\varphi_{d}^{(k'+2)}(x_1,\ldots,x_{k'},v,z)=y$.
\end{itemize}
\item soit $e = \dbllangle 5, k, c\dblrangle$ et les $y+1$ fils portent des
  arbres de calcul attestant $\varphi_{c}^{(k+1)}(0,x_1,\ldots,x_k)=v_0$,
  ..., $\varphi_{c}^{(k+1)}(y,x_1,\ldots,x_k)=v_y$ où $v_y=0$ et
  $v_0,\ldots,v_{y-1} > 0$.
\end{itemize}
\par}

On encode l'arbre $\mathscr{T}$ par l'entier
$\operatorname{code}(\mathscr{T}) := \dbllangle n,
\operatorname{code}(\mathscr{T}_1), \ldots,
\operatorname{code}(\mathscr{T}_s)\dblrangle$ où $n$ est l'étiquette
de la racine et $\mathscr{T}_1,\ldots,\mathscr{T}_s$ les codes des
sous-arbres portés par les $s$ fils de celle-ci.

\end{frame}
%
\begin{frame}
\frametitle{Arbres de calcul $\Rightarrow$ universalité}

Les points-clés :
\begin{itemize}
\item On a $\varphi_e^{(k)}(\underline{x}) = y$ \alert{ssi} il existe
  un arbre de calcul $\mathscr{T}$ l'attestant.
\item Vérifier si $\mathscr{T}$ est un arbre de calcul valable est
  \alert{primitif récursif} en $\operatorname{code}(\mathscr{T})$.
  (On peut vérifier les règles à chaque nœud avec des boucles
  bornées.)
\item De même, extraire $e,\underline{x},y$ de $\mathscr{T}$ est
  primitif récursif.
\end{itemize}

\bigskip

D'où l'algorithme « universel » pour calculer
$\varphi_e^{(k)}(\underline{x})$ en fonction de $e,\underline{x}$ :
\begin{itemize}
\item parcourir $n=0,1,2,3,4,\ldots$ (boucle non bornée),
\item pour chacun, tester s'il code un arbre de calcul valable de
  $\varphi_e^{(k)}(\underline{x})$,
\item si oui, terminer et renvoyer le $y$ contenu.
\end{itemize}

La boucle non-bornée est précisément ce que permet $\mu$.  Tout le
reste est p.r.

$\Rightarrow$ Ceci montre l'existence de $u$ (code de l'algorithme
décrit ci-dessus).

\bigskip

\textcolor{orange}{Ne pas coder un interpréteur comme ça dans la vraie vie !}

\end{frame}
%
\begin{frame}
\label{normal-form-theorem}
\frametitle{Théorème de la forme normale}

On a montré un peu plus que l'universalité : on peut exécuter
n'importe quel algorithme avec une \alert{unique boucle non bornée}.
Plus exactement :

\bigskip

\itempoint\textbf{Théorème de la forme normale} (Kleene) : il existe
un prédicat p.r. $T$ sur $\mathbb{N}^3$ et une fonction p.r. $U \colon
\mathbb{N} \to \mathbb{N}$ tels que :
\[
\varphi_e^{(k)}(x_1,\ldots,x_k)
= U(\mu T(e,\dbllangle x_1,\ldots,x_k\dblrangle))
\]

Précisément, $T(n, e,\dbllangle x_1,\ldots,x_k\dblrangle)$ teste si
$n$ est le code d'un arbre de calcul valable de
$\varphi_e^{(k)}(\underline{x})$, et $U$ extrait le résultat de cet
arbre.

\medskip

\centerline{*}

\medskip

Exemple d'application : \textbf{lancement en parallèle} :
\[
U(\mu(T(e_1,\dbllangle\underline{x}\dblrangle)\text{~ou~}T(e_2,\dbllangle\underline{x}\dblrangle)))
\]
définit (de façon p.r. en $e_1,e_2$) un $e$ tel que
\[
\varphi_e(\underline{x}){\downarrow} \;\Longleftrightarrow\;
\varphi_{e_1}(\underline{x}){\downarrow}\text{~ou~}
\varphi_{e_2}(\underline{x}){\downarrow}
\]

\end{frame}
%
\begin{frame}
\frametitle{Théorème s-m-n (version générale récursive)}

Exactement comme la version p.r. :

\smallskip

\itempoint\textbf{Théorème s-m-n} (Kleene) : il existe $s_{m,n} \colon
\mathbb{N}^{m+1} \to \mathbb{N}$ p.r. telle que
\[
(\forall e,\underline{x},\underline{y})\quad
\varphi^{(n)}_{s_{m,n}(e,x_1,\ldots,x_m)}(y_1,\ldots,y_n) =
\varphi^{(m+n)}_e(x_1,\ldots,x_m,\,y_1,\ldots,y_n)
\]

\bigskip

Noter que $s_{m,n}$ est \alert{primitive récursive} même si on
s'intéresse ici aux fonctions générales récursives.

\medskip

Les manipulations de programmes sont \textcolor{blue}{typiquement
  p.r.} (même si les programmes manipulés sont des fonctions générales
récursives).

\end{frame}
%
\begin{frame}
\frametitle{Arité et encodage des tuples}

{\footnotesize\textcolor{gray}{Remarque qui aurait dû être faite
    avant ?}\par}

\bigskip

Pour tout $k \geq q$, les fonctions
\[
\left\{
\begin{array}{l}
\mathbb{N}^k \to \mathbb{N}\\
(x_1,\ldots,x_k) \mapsto \dbllangle x_1,\ldots,x_k\dblrangle
\end{array}
\right.
\quad\text{~et~}\quad
\left\{
\begin{array}{l}
\mathbb{N} \to \mathbb{N}\\
\dbllangle x_1,\ldots,x_k\dblrangle \mapsto x_i
\end{array}
\right.
\]
sont p.r.  Par conséquent,
\[
f\colon\mathbb{N}^k \dasharrow \mathbb{N}\text{~récursive}
\;\Longleftrightarrow\;
\left\{
\begin{array}{l}
\mathring f\colon\mathbb{N} \dasharrow \mathbb{N}\\
\hphantom{f\colon}
\dbllangle x_1,\ldots,x_k\dblrangle \mapsto f(x_1,\ldots,x_k)
\end{array}
\right.\text{~récursive}
\]
et de plus, un numéro $e$ de $f$ (i.e., $f = \varphi^{(k)}_e$) se
calcule de façon p.r. à partir d'un numéro $e'$ de $\mathring f$
(i.e., $\mathring f = \varphi^{(1)}_{e'}$) et vice versa.

\medskip

Ceci justifie d'omettre parfois abusivement l'arité

(par défaut, « $\varphi_e$ » désigne « $\varphi^{(1)}_e$ »).

\bigskip

{\footnotesize Même chose, \textit{mutatis mutandis} (avec $\psi$)
  pour les fonctions p.r. elles-mêmes.\par}

\end{frame}
%
\begin{frame}
\label{kleene-recursion-theorem}
\frametitle{Le théorème de récursion de Kleene (version générale récursive)}

Exactement comme la version p.r. :

\smallskip

\itempoint\textbf{Théorème} (Kleene) : si $h \colon \mathbb{N}^{k+1}
\dasharrow \mathbb{N}$ est récursive, il existe $e$ tel que
\[
(\forall\underline{x})\quad \varphi^{(k)}_e(\underline{x}) =
h(e,\underline{x})
\]
Plus précisément, il existe $b \colon \mathbb{N}^2 \to \mathbb{N}$
p.r. telle que $e := b(k,d)$ vérifie
\[
(\forall d)\,(\forall\underline{x})\quad \varphi^{(k)}_e(\underline{x}) = \varphi^{(k+1)}_d(e,\underline{x})
\]

\bigskip

\underline{Même preuve :} soit $s := s_{1,k}$ donné par le théorème
s-m-n.  La fonction $(t,\underline{x}) \mapsto
h(s(t,t),\underline{x})$ est générale récursive, disons $=
\varphi_c^{(k+1)}(t,\underline{x})$.  Alors
\[
\varphi_{s(c,c)}^{(k)}(\underline{x})
= \varphi_{c}^{(k+1)}(c, \underline{x})
= h(s(c,c),\underline{x})
\]
donc $e := s(c,c)$ convient.  Les fonctions $d \mapsto c$ et $c \mapsto e$
sont p.r.\qed

\bigskip

\textcolor{blue}{\textbf{Moralité :}} \alert{on peut donner aux programmes accès à leur
  propre numéro} (= « code source »), cela ne change rien.

\end{frame}
%
\begin{frame}
\frametitle{Le théorème du point fixe de Kleene-Rogers}

Reformulation du théorème de récursion utilisant l'universalité :

\smallskip

\itempoint\textbf{Théorème} (Kleene-Rogers) : si $F \colon \mathbb{N}
\dasharrow \mathbb{N}$ est récursive et $k\in\mathbb{N}$, il existe
$e$ tel que
\[
\varphi_e^{(k)} = \varphi_{F(e)}^{(k)}
\]

\bigskip

\underline{Preuve :} $h\colon (e,\underline{x}) \mapsto
\varphi_{F(e)}^{(k)}(\underline{x})$ est récursive car $e \mapsto
F(e)$ l'est et que $(e',\underline{x}) \mapsto
\varphi_{e'}^{(k)}(\underline{x})$ l'est (universalité).  Par le
théorème de récursion, il existe $e$ tel que
$\varphi^{(k)}_e(\underline{x}) = h(e,\underline{x}) =
\varphi_{F(e)}^{(k)}(\underline{x})$.\qed

\bigskip

\textcolor{blue}{\textbf{Moralité :}} quelle que soit la
transformation $F$ calculable effectuée sur le source d'un programme,
il y a un programme $e$ qui fait la même chose que son
transformé $F(e)$.

\end{frame}
%
\begin{frame}
\label{recursion-from-kleene-recursion-theorem}
\frametitle{Récursion !}

Le langage des fonctions générales récursives,
\textcolor{orange}{malgré le nom} ne permet pas les définitions par
appels récursifs.

\smallskip

{\footnotesize Uniquement des opérations élémentaires, appels de
  fonctions précédemment définies, boucles.\par}

\bigskip

Comment permettre quand même les appels récursifs ?

\smallskip

\alert{Par le théorème de récursion de Kleene !} (ou théorème du point fixe) :

\begin{itemize}
\item je veux définir (comme fonction générale récursive) une fonction
  $f$ dont la définition fait appel à $f$ elle-même :
\item par le théorème de récursion de Kleene (« astuce de Quine »), je
  peux supposer que $f$ a accès à son propre numéro (« code source »),
\item je convertis chaque appel à $f$ depuis $f$ en un appel à la
  fonction universelle (interpréteur) sur le numéro de $f$.
\end{itemize}

\bigskip

\textcolor{orange}{Ne pas implémenter la récursion comme ça dans la vraie vie !}

\end{frame}
%
\begin{frame}
\frametitle{\textit{Kids, don't try this at home !}}

Pseudocode :

\smallskip

{\footnotesize\texttt{%
fibonacci(n) \{\newline
str = "self = \textbackslash"fibonacci(n) \{\textbackslash \textbackslash nstr = \textbackslash" + quote(str) + str;\textbackslash n\textbackslash\newline
if (n==0 || n==1) return n;\textbackslash n\textbackslash\newline
return interpret(self, n-1) + interpret(self, n-2);\textbackslash n\textbackslash\newline
\}";\newline
self = "fibonacci(n) \{\textbackslash nstr = " + quote(str) + str;\newline
if (n==0 || n==1) return n;\newline
return interpret(self, n-1) + interpret(self, n-2);\newline
\}
}\par}

\medskip

\centerline{*}

\medskip

\textbf{Défi :} trouver explicitement un $e$ tel que $\varphi^{(3)}_e$
soit la fonction d'Ackermann.

\smallskip

(La fonction d'Ackermann a été définie par des appels récursifs donc
elle est bien censée être calculable.)

\end{frame}
%
\begin{frame}
\frametitle{Le problème de l'arrêt}

{\footnotesize Le terme « problème de l'arrêt » prendra plus de sens
  pour les machines de Turing.\par}

\medskip

\itempoint\textbf{Problème :} donné un programme $e$ (mettons d'arité
$k=1$) et une entrée $x$ à ce programme, comment savoir si
l'algorithme $e$ termine (c'est-à-dire $\varphi^{(1)}_e(x)\downarrow$)
ou non ($\varphi^{(1)}_e(x)\uparrow$) sur cette entrée ?

\medskip

Cette question est-elle \alert{algorithmique} ?

\bigskip

\textbf{Réponse} de Turing : \alert{non}.

\bigskip

\itempoint \textcolor{blue}{Intuition de la preuve :} supposons que
j'aie un moyen algorithmique pour savoir si un algorithme termine ou
pas, je peux lui demander ce que « je » vais faire (astuce de
Quine !), et faire le contraire, ce qui conduit à un paradoxe.

\end{frame}
%
\begin{frame}
\label{undecidability-halting-problem}
\frametitle{L'indécidabilité du problème de l'arrêt}

\itempoint\textbf{Théorème} ([essentiellement] Turing) : il n'existe
pas de fonction récursive $h\colon \mathbb{N}^2 \to \mathbb{N}$ telle
que $h(e,x) = 1$ si $\varphi^{(1)}_e(x)\downarrow$ et $h(e,x) = 0$ si
$\varphi^{(1)}_e(x)\uparrow$.

\bigskip

\underline{Preuve :} par l'absurde : si un tel $h$ existe, alors la
fonction
\[
v\colon (e,x) \mapsto \left\{
\begin{array}{ll}
42&\text{~si~}h(e,x) = 0\\
\uparrow&\text{~si~}h(e,x) = 1\\
\end{array}
\right.
\]
est générale récursive (tester is $h(e,x)=0$, si oui renvoyer $42$,
sinon faire une boucle infinie, p.ex. $\mu(x\mapsto 1)$).

\medskip

Par le théorème de récursion de Kleene, il existe $e$ tel que
$\varphi^{(1)}_e(x) = v(e,x)$.

\medskip

Si $\varphi^{(1)}_e(x)\downarrow$ alors $h(e,x) = 1$ donc
$v(e,x)\uparrow$ donc $\varphi^{(1)}_e(x)\uparrow$, contradiction.

Si $\varphi^{(1)}_e(x)\uparrow$ alors $h(e,x) = 0$ donc
$v(e,x)\downarrow$ donc $\varphi^{(1)}_e(x)\downarrow$,
contradiction.\qed

\end{frame}
%
\begin{frame}
\label{undecidability-halting-problem-redux}
\frametitle{L'indécidabilité du problème de l'arrêt : redite}

{\footnotesize Notons $\varphi$ pour $\varphi^{(1)}$.\par}

\smallskip

\itempoint\textbf{Théorème} ([essentiellement] Turing) : il n'existe
pas de fonction récursive $h\colon \mathbb{N}^2 \to \mathbb{N}$ telle
que $h(e,x) = 1$ si $\varphi_e(x)\downarrow$ et $h(e,x) = 0$ si
$\varphi_e(x)\uparrow$.

\bigskip

\underline{Preuve} (incluant celle du théorème de récursion) :
considérons la fonction $v\colon \mathbb{N} \dasharrow \mathbb{N}$ qui
à $e$ associe $42$ si $h(e,e)=0$ et $\uparrow$ (non définie) si
$h(e,e)=1$.

Supposons par l'absurde $h$ est calculable : alors cette fonction
(partielle) $v$ est calculable, disons $v = \varphi_c$.

Si $\varphi_c(c)\downarrow$ alors $h(c,c)=1$ donc $v(c)\uparrow$,
c'est-à-dire $\varphi_c(c)\uparrow$, contradiction.\newline Si
$\varphi_c(c)\uparrow$ alors $h(c,c)=0$ donc $v(c)\downarrow$,
c'est-à-dire $\varphi_c(c)\downarrow$, contradiction.\qed

\bigskip

C'est un \alert{argument diagonal} : on utilise $h$ pour construire
une fonction qui diffère en tout point de la diagonale $c \mapsto
\varphi_c(c)$, donc elle ne peut pas être une $\varphi_c$.

\medskip

{\footnotesize Pour les fonctions p.r. (qui terminent toujours !), le
  même argument diagonal donnait l'inexistence d'un programme
  universel (transp. \ref{primitive-recursive-no-universality}).\par}

\end{frame}
%
\begin{frame}
\frametitle{Comparaison fonctions primitives récursives et générales récursives}

\textcolor{violet}{Récapitulation :}

\medskip

\itempoint Les fonctions p.r. sont totales ; les générales récursives
sont possiblement partielles.

\medskip

\itempoint Les fonctions p.r. sont un langage limité (pas de boucle
non bornées a priori) ; les générales récursives coïncideront avec les
fonctions « calculables » (équivalence avec machines de Turing et
$\lambda$-calcul à voir).

\medskip

\itempoint Les fonctions p.r. ne permettent pas d'interpréter les
fonctions p.r. ; les générales récursives peuvent s'interpréter
elles-mêmes (universalité) et donc réaliser n'importe quelle sorte
d'appels récursifs.

\medskip

\itempoint Le problème de l'arrêt pour les fonctions p.r. est trivial
(elles sont totales !) ; pour les fonctions générales récursives, il
est indécidable (= pas calculable par une fonction générale
récursive).

\end{frame}
%
\section{Machines de Turing}
\begin{frame}
\frametitle{Machines de Turing : explication informelle}

La \textbf{machine de Turing} est une modélisation d'un ordinateur
extrêmement simple, réalisant des calculs indiscutablement finitistes.

\medskip

C'est une sorte d'automate doté d'un \textbf{état} interne pouvant
prendre un nombre fini de valeurs, et d'une mémoire illimitée sous
forme de \textbf{bande} linéaire divisée en cellules (indéfiniment
réécrivibles), chaque cellule pouvant contenir un \textbf{symbole}.

\medskip

La machine peut observer, outre son état interne, une unique case de
la bande, là où se trouve sa \textbf{tête de lecture/écriture}.

\medskip

Le \textbf{programme} de la machine indique, pour chaque combinaison
de l'état interne et du symbole lu par la tête :
\begin{itemize}
\item dans quel état passer,
\item quel symbole écrire à la place de la tête,
\item la direction dans laquelle déplacer la tête (gauche ou droite).
\end{itemize}

\medskip

La machine suit son programme jusqu'à tomber dans un état spécial $0$
(« arrêt »).

\end{frame}
%
\begin{frame}
\frametitle{Machines de Turing : définition}

Une \textbf{machine de Turing} (déterministe) à ($1$ bande,
$2$ symboles et) $m\geq 2$ états est la donnée de :
\begin{itemize}
\item un ensemble fini $Q$ de cardinal $m$ d'\textbf{états}, qu'on
  identifiera à $\{0,\ldots,m-1\}$,
\item un ensemble $\Sigma$ de (ici) $2$ \textbf{symboles de bande}
  qu'on identifiera à $\{0,1\}$,
\item une fonction
\[
\delta \colon (Q\setminus\{0\}) \times \Sigma \to Q \times \Sigma \times \{\texttt{L},\texttt{R}\}
\]
appelé \textbf{programme} de la machine.
\end{itemize}

{\footnotesize (Il y a donc $(4m)^{2(m-1)}$ machines à $m$ états.)\par}

\bigskip

Une \textbf{configuration} d'une telle machine est la donnée de :
\begin{itemize}
\item un élément $q \in Q$ appelé l'\textbf{état courant},
\item une fonction $\beta\colon \mathbb{Z} \to \Sigma$ ne prenant
  qu'\alert{un nombre fini} de valeurs $\neq 0$, appelée la
  \textbf{bande},
\item un entier $i \in \mathbb{Z}$ appelé la \textbf{position de la
  tête} sur la bande.
\end{itemize}

\end{frame}
%
\begin{frame}
\frametitle{Machines de Turing : exécution d'une étape}

Si $(q,\beta,i)$ est une configuration de la machine de Turing où
$q\neq 0$, et $\delta$ le programme, la \textbf{configuration
  suivante} est $(q',\beta',i')$ où :
\begin{itemize}
\item $(q',y,d) = \delta(q,\beta(i))$ est l'\textbf{instruction
  exécutée},
\item $q'$ est le \textbf{nouvel état},
\item $i' = i-1$ si $d=\texttt{L}$ et $i' = i+1$ si $d=\texttt{R}$,
\item $\beta'(j) = \beta(j)$ pour $j\neq i$ tandis que $\beta'(i) = y$.
\end{itemize}

\bigskip

\emph{En clair :} le programme indique, pour chaque configuration d'un
état $\neq 0$ et d'un symbole $x = \beta(i)$ lu sur la bande :
\begin{itemize}
\item le nouvel état $q'$ dans lequel passer,
\item le symbole $y$ à écrire à la place de $x$ à l'emplacement $i$ de
  la bande,
\item la direction dans laquelle déplacer la tête (gauche ou droite).
\end{itemize}

\end{frame}
%
\begin{frame}
\frametitle{Machines de Turing : exécution complète}

\itempoint Si $C = (q,\beta,i)$ est une configuration d'une machine de
Turing, la \textbf{trace d'exécution} à partir de $C$ est la suite
finie ou infinie $C^{(0)},C^{(1)},C^{(2)},\ldots$, où
\begin{itemize}
\item $C^{(0)} = C$ est la configuration donnée (configuration initiale),
\item si $C^{(n)} = (q^{(n)},\beta^{(n)},i^{(n)})$ avec $q^{(n)}=0$
  alors la suite s'arrête ici, on dit que \textbf{la machine
    s'arrête}, que $C^{(n)}$ est la \textbf{configuration finale}, et
  que l'exécution a duré $n$ \textbf{étapes},
\item sinon, $C^{(n+1)}$ est la configuration suivante (définie
  transp. précédent).
\end{itemize}

\bigskip

\emph{En clair :} la machine continue à exécuter des instructions tant
qu'elle n'est pas tombée dans l'état $0$.  Elle s'arrête quand elle
tombe dans l'état $0$.

\end{frame}
%
\begin{frame}
\label{simulation-of-turing-machines-by-recursive-functions}
\frametitle{Simulation des machines de Turing par les fonctions récursives}

\itempoint On peut coder un programme et/ou une configuration sous
forme d'entiers naturels.

{\footnotesize Le ruban a un nombre \alert{fini} de symboles $\neq 0$,
  donc on peut le coder par la liste de leurs positions comptées,
  disons, à partir du symbole $\neq 0$ le plus à gauche.\par}

\bigskip

\itempoint La fonction $(M,C) \mapsto C'$ qui à une machine de Turing
$M$ et une configuration $C$ de $M$ associe la configuration suivante
\textbf{est p.r.}

\medskip

\itempoint Conséquence : la fonction $(n,M,C) \mapsto C^{(n)}$ qui à
$n\in\mathbb{N}$ et une machine de Turing $M$ et une configuration $C$
de $M$ associe la configuration atteinte après $n$ étapes d'exécution,
\textbf{est p.r.}

\medskip

\itempoint La fonction qui à $(M,C)$ associe la configuration finale
(et/ou le nombre d'étapes d'exécution) \alert{si la machine s'arrête},
et $\uparrow$ (non définie) si elle ne s'arrête pas, est
\textbf{générale récursive}.

\bigskip

\textcolor{blue}{\textbf{Moralité :}} les fonctions récursives peuvent
simuler les machines de Turing.

\end{frame}
%
\begin{frame}
\frametitle{Calculs sur machines de Turing : une convention}

On dira qu'une fonction $f\colon \mathbb{N}^k \dasharrow \mathbb{N}$
est \textbf{calculable par machine de Turing} lorsqu'il existe une
machine de Turing qui, pour tous $x_1,\ldots,x_k$ :
\begin{itemize}
\item part de la configuration initiale suivante : l'état est $1$, les
  symboles $\beta(j)$ du ruban pour $j<0$ sont arbitraires (tous $0$
  sauf un nombre fini), la tête est à l'emplacement $0$,
\item les symboles $\beta(j)$ pour $j\geq 0$ du ruban initial forment
  le mot suivant :
\[
0 1^{x_1} 0 1^{x_2} 0 \cdots 0 1^{x_k} 0
\]
(suivi d'une infinité de $0$), c'est-à-dire $\beta(0)=0$, $\beta(j)=1$
si $1\leq j\leq x_1$, $\beta(1+x_1)=0$, $\beta(j)=1$ si $2+x_1\leq
j\leq 1+x_1+x_2$, etc.,
\item si $f(x_1,\ldots,x_k)\uparrow$, la machine ne s'arrête pas,
\item si $f(x_1,\ldots,x_k){\downarrow} = y$, la machine s'arrête avec
  la tête à l'emplacement $0$ (le même qu'au départ), le ruban
  $\beta(j)$ non modifié pour $j<0$, et
\item les symboles $\beta(j)$ pour $j\geq 0$ du ruban final forment le
  mot $0 1^y 0$ (suivi d'une infinité de $0$) {\footnotesize (« codage
    unaire » de $y$)}.
\end{itemize}

\end{frame}
%
\begin{frame}
\frametitle{Calculs par les machines de Turing des fonctions récursives}

\itempoint On peut montrer par induction suivant la
déf\textsuperscript{n} de $\mathbf{R}$ que \alert{toute fonction
  générale récursive est calculable par machine de Turing} avec les
conventions du transp. précédent.

\bigskip

\itempoint La démonstration est fastidieuse mais pas difficile : il
s'agit essentiellement de programmer en machine de Turing chacune des
formes de construction des fonctions générales récursives
(projections, constantes, successeur, composition, récursion
primitive, $\mu$-récursion).

\bigskip

\itempoint Les conventions faites, notamment le fait d'ignorer et de
ne pas modifier $\beta(j)$ pour $j<0$, permettent à l'induction dans
la preuve de fonctionner.

\smallskip

{\footnotesize Par exemple, pour la composition, on va utiliser cette
  propriété pour « sauvegarder » les $x_1,\ldots,x_k$ initiaux, ainsi
  que les valeurs de $g_j(\underline{x})$ calculées, lorsqu'on appelle
  chacune des fonctions $g_1,\ldots,g_\ell$ (à chaque fois, on les
  recopie $x_1,\ldots,x_k$ à droite des valeurs à ne pas toucher, et
  on appelle la machine calculant $g_j$ sur ces valeurs
  recopiées).\par}

\end{frame}
%
\begin{frame}
\frametitle{Équivalence entre machines de Turing et fonctions récursives}

\itempoint Toute fonction générale récursive $\mathbb{N}^k \dasharrow
\mathbb{N}$ est calculable par machine de Turing (sous les conventions
données) : transp. précédent.

\bigskip

\itempoint Réciproquement, toute fonction $\mathbb{N}^k \dasharrow
\mathbb{N}$ calculable par machine de Turing sous ces conventions est
générale récursive, car les fonctions récursives peuvent simuler les
machines de Turing, calculer une configuration initiale convenable, et
décoder la configuration finale
(cf. transp. \ref{simulation-of-turing-machines-by-recursive-functions}).

\bigskip

\itempoint Bref, $f\colon \mathbb{N}^k \dasharrow \mathbb{N}$ est
calculable par machine de Turing \alert{ssi} elle est générale
récursive.

\bigskip

\itempoint De plus, cette équivalence est \alert{constructive} : il
existe des fonctions p.r. :
\begin{itemize}
\item l'une prend en entrée le numéro $e$ d'une fonction générale
  récursive (et l'arité $k$) et renvoie le code d'une machine de
  Turing qui calcule cette $\varphi_e^{(k)}$,
\item l'autre prend en entrée le code d'une machine de Turing qui
  calcule une fonction $f$, et son arité $k$, et renvoie un numéro $e$
  de $f$ dans les fonctions générales récursives $f =
  \varphi_e^{(k)}$.
\end{itemize}

\end{frame}
%
\begin{frame}
\frametitle{Machines de Turing : variations}

On a choisi ici une notion de machine de Turing assez restreinte
($1$ bande, $2$ symboles de bande).  Il existe toutes sortes de
variations :
\begin{itemize}
\item machines à plusieurs bandes (mais en nombre fini ; le programme
  choisit en fonction du symbole lu sur chaque bande, et écrit et
  déplace chaque tête indépendamment), voire à plusieurs têtes par
  bande, parfois avec des bandes en lecture seule (pour les entrées),
  ou en écriture seule (pour les sorties),
\item autres symboles que $0$ et $1$ (mais en nombre fini),
\item machine non-déterministe (plusieurs instructions possibles dans
  une configuration donnée ; la machine termine si au moins l'un des
  chemins d'exécution termine).
\end{itemize}

\bigskip

Du point de vue \alert{calculabilité}, ces modifications ne rendent
pas la machine plus puissante, et, sauf, cas dégénérés (p.ex., un seul
symbole sur le ruban !) elles ne la rendent pas moins puissante non
plus.  Ceci confirme la robustesse du modèle de Church-Turing.

\smallskip

{\footnotesize Pour la \alert{complexité}, en revanche, c'est une
  autre affaire.\par}

\end{frame}
%
\begin{frame}
\frametitle{Machines de Turing : reprise de résultats déjà vus}

\itempoint\textbf{Universalité :} pour un codage raisonnable, il
existe une machine de Turing « universelle » qui prend en entrée sur
sa bande le programme d'une autre machine de Turing $M$, et une
configuration initiale $C$ pour celle-ci, et qui simule l'exécution de
$M$ sur $C$ (notamment, elle s'arrête ssi $M$ s'arrête).

\bigskip

\itempoint\textbf{Forme normale :} la fonction $(n,M,C) \mapsto
C^{(n)}$ qui à $n\in\mathbb{N}$ et une machine de Turing $M$ et une
configuration $C$ de $M$ associe la configuration après $n$ étapes
d'exécution, est p.r., et en particulier, calculable par une machine
de Turing.

\smallskip

$\Rightarrow$ En particulier, on peut tester algorithmiquement si une
machine de Turing donnée, depuis une configuration initiale donnée,
s'arrête \emph{en moins de $n$ étapes}.

\bigskip

\itempoint\textbf{Indécidabilité du problème de l'arrêt :} la fonction
qui à $(M,C)$ associe $1$ si la machine de Turing s'arrête en partant
de la configuration initiale $C$, et $0$ sinon, \alert{n'est pas
  calculable}.

\smallskip

$\Rightarrow$ On ne peut pas tester algorithmiquement si une machine de
Turing donnée, depuis une configuration initiale donnée, s'arrête « un
jour ».

\end{frame}
%
\begin{frame}
\label{undecidability-halting-problem-turing-machines-pristine-start}
\frametitle{Indécidabilité du problème de l'arrêt (départ bande vierge)}

On ne peut même pas tester algorithmiquement si une machine de Turing
s'arrête à partir d'une bande vierge :

\smallskip

\itempoint\textbf{Indécidabilité du problème de l'arrêt :} la fonction
qui à $M$ associe $1$ si la machine de Turing s'arrête en partant de
la configuration vierge $C_0$ (c'est-à-dire celle où $\beta = 0$, état
initial $1$), et $0$ sinon, \alert{n'est pas calculable}.

\medskip

\underline{Preuve :} Supposons par l'absurde qu'on puisse tester
algorithmiquement si une machine de Turing s'arrête à partir d'une
configuration vierge.  On va montrer qu'on peut tester si une machine
de Turing $M$ s'arrête à partir de $C$ quelconque.

\smallskip

Données $M$ et $C$, on peut algorithmiquement calculer une machine $N$
qui « prépare » $C$ à partir de la configuration vierge $C_0$, donc
une machine $M^*$ qui exécute successivement $N$ puis $M$
\textcolor{teal}{($\leftarrow$ ceci est un théorème s-m-n)}.

\smallskip

Ainsi $M^*$ (calculé algorithmiquement) termine sur $C_0$ ssi $M$
termine sur $C$.

\smallskip

Donc tester la terminaison de $M^*$ permettrait de tester celle de
$M$ sur $C$, ce qui n'est pas possible
\textcolor{teal}{($\leftarrow$ ceci est une preuve « par
  réduction »)}.\qed

\end{frame}
%
\begin{frame}
\frametitle{Le castor affairé}

\itempoint La fonction \textbf{castor affairé} associe à $m$ le nombre
maximal $B(m)$ d'étapes d'exécution d'une machine de Turing à $\leq m$
états \alert{qui termine} (à partir de la configuration vierge $C_0$).

\medskip

\itempoint La fonction $B$ croît \alert{trop vite pour être
  calculable} :
\[
\forall f\colon\mathbb{N}\to\mathbb{N}\text{~calculable}.\quad
\exists m\in\mathbb{N}.\quad
(B(m) > f(m))
\]

\medskip

\underline{Preuve :} supposons au contraire $\forall m\in\mathbb{N}.\;
(B(m) \leq f(m))$ avec $f\colon\mathbb{N}\to\mathbb{N}$ calculable.
Donnée une machine de Turing $M$, on peut alors algorithmiquement
décider si $M$ s'arrête à partir de $C_0$ :
\begin{itemize}
\item calculer $f(m)$ où $m$ est le nombre d'états de $M$,
\item exécuter $M$ à partir de $C_0$ pendant $f(m)$ étapes (ce nombre
  est $\geq B(m)$ par hypothèse),
\item si elle a terminé en temps imparti, $M$ termine sur $C_0$, et on
  renvoie « oui » ; sinon, elle ne termine jamais par définition de
  $B(m)$, on renvoie « non ».
\end{itemize}
Ceci est impossible donc $f$ n'est pas calculable.\qed

\end{frame}
%
\begin{frame}
\frametitle{Le castor affairé (amélioration)}

{\footnotesize $B(m)=$ nombre maximal d'étapes d'exécution d'une
  machine de Turing à $\leq m$ états \alert{qui termine} à partir
  d'une bande vierge.\par}

\medskip

\itempoint On peut faire mieux : $B$ \alert{domine} toute fonction
calculable :
\[
\forall f\colon\mathbb{N}\to\mathbb{N}\text{~calculable}.\quad
\exists m_0\in\mathbb{N}.\quad
\forall m\geq m_0.\quad
(B(m) > f(m))
\]

\medskip

{\footnotesize

\underline{Preuve :} soit $f\colon\mathbb{N}\to\mathbb{N}$ calculable.
Soit $\gamma(r) = A(2,r,r)$ (en fait, $2^r$ doit suffire ; noter
$\gamma$ croissante).  Pour $r \in \mathbb{N}$, on considère la
machine de Turing $M_r$ qui
\begin{itemize}
\item prépare $r$, calcule $\gamma(r+1)$ puis $f(0) + f(1) + \cdots +
  f(\gamma(r+1)) + 1$,
\item attend ce nombre-là d'étapes, et termine.
\end{itemize}
Le nombre d'états de $M_r$ est une fonction p.r. $b(r)$ de $r$ (même
$b(r) = r + \mathrm{const}$ convient).  Pour $r\geq r_0$ on a $b(r)
\leq \gamma(r)$.  Soit $m_0 = \gamma(r_0)$.  Si $m \geq m_0$, soit
$r\geq r_0$ tel que $\gamma(r) \leq m \leq \gamma(r+1)$.  Alors $M_r$
calcule $\cdots+f(m)+\cdots+1$, donc attend $>f(m)$ étapes.  Donc
$B(b(r)) > f(m)$.  Mais $b(r) \leq \gamma(r) \leq m$ donc $B(m) >
f(m)$.\qed

\par}

\medskip

\centerline{*}

\medskip

\itempoint Variations du castor affairé : nombre de symboles écrits
sur la bande, $n \mapsto \max\{\varphi_e(e) : 0\leq e\leq
n\text{~et~}\varphi_e(e)\downarrow\}$ (mêmes propriétés).

\end{frame}
%
\section{Décidabilité et semi-décidabilité}
\begin{frame}
\frametitle{Terminologie calculable/décidable}

\itempoint Une fonction partielle $f\colon \mathbb{N}^k \dasharrow
\mathbb{N}$ est dite \textbf{calculable} (partielle) lorsqu'elle est
(c'est équivalent) :
\begin{itemize}
\item générale récursive,
\item calculable par machine de Turing,
\item \textcolor{brown}{à voir $\rightarrow$} représentable dans le $\lambda$-calcul.
\end{itemize}

\bigskip

\itempoint Une partie $A \subseteq \mathbb{N}^k$ est dite
\textbf{décidable} lorsque sa fonction indicatrice
$\mathbb{N}^k\to\mathbb{N}$
\[
\mathbf{1}_A\colon n \mapsto \left\{
\begin{array}{ll}
1&\text{~si~}n\in A\\
0&\text{~si~}n\not\in A\\
\end{array}
\right.
\]
est calculable (répondre « oui » ou « non » selon que $n\in A$ ou $n\not\in A$).

\bigskip

\itempoint Une partie $A \subseteq \mathbb{N}^k$ est dite
\textbf{semi-décidable} lorsque sa fonction partielle « semi-indicatrice »
$\mathbb{N}\dasharrow\mathbb{N}$ (d'ensemble de définition $A$)
\[
n \mapsto \left\{
\begin{array}{ll}
1&\text{~si~}n\in A\\
\uparrow&\text{~si~}n\not\in A\\
\end{array}
\right.
\]
est calculable (répondre « oui » ou « ... » selon que $n\in A$ ou $n\not\in A$).

\end{frame}
%
\begin{frame}
\frametitle{Fluctuations terminologiques}

\itempoint Synonymes de \textbf{calculable} pour une fonction
partielle $\mathbb{N}^k \dasharrow \mathbb{N}$ :
\begin{itemize}
\item « semi-calculable » (réservant « calculable » pour les fonctions \emph{totales}),
\item « (générale) récursive ».
\end{itemize}

\bigskip

\itempoint Synonymes de \textbf{décidable} pour une partie $\subseteq
\mathbb{N}^k$ :
\begin{itemize}
\item « calculable »,
\item « récursive ».
\end{itemize}

\bigskip

\itempoint Synonymes de \textbf{semi-décidable} pour une partie $\subseteq
\mathbb{N}^k$ :
\begin{itemize}
\item « semi-calculable »,
\item « calculablement énumérable »,
\item « récursivement énumérable ».
\end{itemize}
{\footnotesize (La raison du mot « énumérable » sera expliquée après.)\par}

\end{frame}
%
\begin{frame}
\frametitle{Décidable = semi-décidable de complémentaire semi-décidable}

\itempoint Si $A \subseteq \mathbb{N}^k$ est décidable, alors son
complémentaire $\complement A := \mathbb{N}^k \setminus A$ l'est
aussi.

{\footnotesize \underline{Preuve :} échanger $0$ et $1$ dans la
  réponse. \qedsymbol\par}

\medskip

\itempoint Si $A$ est décidable, alors $A$ est semi-décidable.

{\footnotesize \underline{Preuve :} si réponse $0$, faire une boucle
  infinie. \qedsymbol\par}

\medskip

\itempoint Donc : si $A$ est décidable, alors $A$ et $\complement A$
sont semi-décidables.

\bigskip

\itempoint La réciproque est également valable : si $A$ et
$\complement A$ sont semi-décidables alors $A$ est décidable.

\medskip

\textcolor{blue}{Idée :} lancer « en parallèle » un algorithme qui
semi-décide $A$ et un qui semi-décide $\complement A$ ; l'un des deux
finira par donner la réponse voulue.

\medskip

\textcolor{brown}{Mais que signifie « lancer en parallèle » ici ?}

\end{frame}
%
\begin{frame}
\frametitle{Lancement en parallèle}

On suppose que :
\begin{itemize}
\item $\varphi_{e_1}(\underline{x})\downarrow$ ssi $\underline{x} \in A$
\item $\varphi_{e_2}(\underline{x})\downarrow$ ssi $\underline{x} \not\in A$
\end{itemize}
Comment décider si $\underline{x} \in A$ en terminant à coup sûr ?

\bigskip

Grâce au \alert{th. de la forme normale}
(transp. \ref{normal-form-theorem}) : il y a un prédicat $T$ p.r. tel
que
\begin{itemize}
\item $\varphi_{e_1}(\underline{x})\downarrow$ ssi $\exists
  n\in\mathbb{N}.\; T(n,e_1,\dbllangle\underline{x}\dblrangle)$
\item $\varphi_{e_2}(\underline{x})\downarrow$ ssi $\exists
  n\in\mathbb{N}.\; T(n,e_2,\dbllangle\underline{x}\dblrangle)$
\end{itemize}
On a alors $\exists n\in\mathbb{N}.\;
(T(n,e_1,\dbllangle\underline{x}\dblrangle) \text{~ou~}
T(n,e_2,\dbllangle\underline{x}\dblrangle))$ à coup sûr.

\bigskip

Algorithme (terminant à coup sûr) :
\begin{itemize}
\item parcourir $n=0,1,2,3,4,\ldots$ (boucle non bornée),
\item pour chacun, tester si
  $T(n,e_1,\dbllangle\underline{x}\dblrangle)$ et si
  $T(n,e_2,\dbllangle\underline{x}\dblrangle)$,
\item si le premier vaut, renvoyer « oui, $\underline{x}\in A$ », si
  le second vaut, renvoyer « non, $\underline{x}\not\in A$ » (sinon,
  continuer la boucle).
\end{itemize}

\end{frame}
%
\begin{frame}
\frametitle{Lancement en parallèle (variante machines de Turing)}

On suppose que :
\begin{itemize}
\item la machine $M_1$ s'arrête sur $\underline{x}$ ssi $\underline{x}
  \in A$
\item la machine $M_2$ s'arrête sur $\underline{x}$ ssi $\underline{x}
  \not\in A$
\end{itemize}
Comment décider si $\underline{x} \in A$ en s'arrêtant à coup sûr ?

\bigskip

On va simuler $M_1$ et $M_2$ pour $n$ étapes jusqu'à ce que l'une
d'elles s'arrête.

\bigskip

Algorithme (terminant à coup sûr) :
\begin{itemize}
\item parcourir $n=0,1,2,3,4,\ldots$ (boucle non bornée),
\item pour chacun, tester si l'exécution de $M_1$ s'arrête sur
  $\underline{x}$ en $\leq n$ étapes et si l'exécution de $M_2$
  s'arrête sur $\underline{x}$ en $\leq n$ étapes,
\item si le premier vaut, renvoyer « oui, $\underline{x}\in A$ », si
  le second vaut, renvoyer « non, $\underline{x}\not\in A$ » (sinon,
  continuer la boucle).
\end{itemize}

\bigskip

{\footnotesize C'est \alert{exactement la même chose} que dans le
  transp. précédent, avec un nombre d'étapes d'exécution $n$ au lieu
  d'un arbre de calcul (détail sans importance).\par}

\end{frame}
%
\begin{frame}
\frametitle{Problème de l'arrêt}

Le \textbf{problème de l'arrêt} est :
\[
\mathscr{H} := \{(e,x)\in\mathbb{N}^2 : \varphi_e(x)\downarrow\}
\]

\smallskip

\itempoint Il \alert{n'est pas décidable}
(transp. \ref{undecidability-halting-problem}).

\smallskip

\itempoint Il \alert{est} semi-décidable (par universalité : donné
$(e,x)$, on peut exécuter $\varphi_e(x)$, et, s'il termine, renvoyer
« oui »).

\smallskip

\itempoint Donc $\complement\mathscr{H}$ n'est pas semi-décidable.

\bigskip

{\footnotesize

\itempoint Toutes sortes de variantes possibles, p.ex. :
\begin{itemize}
\item $\{e\in \mathbb{N} : \varphi_e(e)\downarrow\}$ n'est pas décidable
  (preuve dans transp. \ref{undecidability-halting-problem-redux})
\item $\{e\in \mathbb{N} : \varphi_e(0)\downarrow\}$ n'est pas décidable
  (théorème s-m-n : $\varphi_e(x) = \varphi_{s(e,x)}(0)$ avec $s$
  p.r. ;
  cf. transp. \ref{undecidability-halting-problem-turing-machines-pristine-start})
\end{itemize}

\par}

\end{frame}
%
\begin{frame}
\frametitle{Image d'un ensemble décidable}

\itempoint Si $A \subseteq \mathbb{N}$ est décidable et $f \colon
\mathbb{N} \to \mathbb{N}$ (totale) calculable, alors l'image
\[
f(A) := \{f(i) : i\in A\}
\]
est semi-décidable.

\smallskip

{\footnotesize \underline{Preuve :} donné $m\in\mathbb{N}$, pour
  semi-décider si $m \in f(A)$, parcourir $i=0,1,2,3\ldots$, et pour
  chacun, décider si $i\in A$ et, si oui, calculer $f(i)$ et comparer
  à $m$.  Si $i\in A$ et $f(i)=m$, renvoyer « oui » ; sinon, continuer
  la boucle.\qed\par}

\bigskip

\itempoint Réciproquement, si $B \subseteq \mathbb{N}$ est
semi-décidable, il existe $A \subseteq \mathbb{N}$ décidable et $f
\colon \mathbb{N} \to \mathbb{N}$ (totale) calculable tels que $B =
f(A)$.

\smallskip

{\footnotesize \underline{Preuve :} soit $e$ tel que $B = \{m :
  \varphi_e(m)\downarrow\}$ ; soit $A$ l'ensemble des $\langle
  n,m\rangle$ tels que $T(n,e,\dbllangle m\dblrangle)$ : alors $A$ est
  décidable et son image par $\langle n,m\rangle \mapsto m$
  est $B$.\qed\par}

{\footnotesize \underline{Redite :} soit $M$ une machine de Turing qui
  s'arrête sur $m$ ssi $m \in B$ ; soit $A$ l'ensemble des $\langle
  n,m\rangle$ tels que $M$ s'arrête sur $m$ en $\leq n$ étapes : alors
  $A$ est décidable et son image par $\langle n,m\rangle \mapsto m$
  est $B$.\qed\par}

\bigskip

\itempoint Variante : $B \subseteq \mathbb{N}$ \emph{non vide} est
semi-décidable ssi il existe $f\colon \mathbb{N} \to \mathbb{N}$
totale calculable telle que $f(\mathbb{N}) = B$.
\textcolor{teal}{D'où le terme « calculablement énumérable ».}

\end{frame}
%
\begin{frame}
\frametitle{Stabilités par opérations booléennes}

Les ensembles \textbf{décidables} sont stables par :
\begin{itemize}
\item réunions finies,
\item intersections finies,
\item complémentaire,
\item \alert{mais pas par} projection $\mathbb{N}^k \to
  \mathbb{N}^{k'}$ (où $k'\leq k$).\\ {\footnotesize (Le problème de
    l'arrêt est une projection d'un ensemble décidable,
    cf. transp. précédent.)}
\end{itemize}

\bigskip

Les ensembles \textbf{semi-décidables} sont stables par :
\begin{itemize}
\item réunions finies (par lancement en parallèle !),
\item intersections finies,
\item projection $\mathbb{N}^k \to \mathbb{N}^{k'}$ (où $k'\leq
  k$),\\ {\footnotesize (Les ensembles semi-décidables sont
    projections d'ensembles décidables donc sont eux-mêmes stables par
    projections, cf. transp. précédent et idées proches.)}
\item \alert{mais pas par complémentaire}.
\end{itemize}

\end{frame}
%
\begin{frame}
\frametitle{Le théorème de Rice : énoncé}

Soit $\textbf{R}^{(1)}$ l'ensemble des fonctions partielles
$\mathbb{N} \dasharrow \mathbb{N}$ calculables (= générales
récursives), et $\Phi \colon e \mapsto \varphi^{(1)}_e$ qui définit
une surjection $\mathbb{N} \to \textbf{R}^{(1)}$.

\medskip

{\footnotesize Si $e$ est l'« intention » (l'algorithme, le
  programme), alors $\Phi(e)$ est l'« extension » (la fonction, i.e.,
  son graphe) définie par $e$.\par}

\bigskip

\itempoint\textbf{Théorème} (Rice) : si $F \subseteq
\textbf{R}^{(1)}$ est un ensemble de fonctions partielles tel que
$\Phi^{-1}(F) := \{e \in \mathbb{N} : \varphi^{(1)}_e \in F\}$ est
\emph{décidable}, alors $F = \varnothing$ ou $F = \textbf{R}^{(1)}$.

\bigskip

\textcolor{blue}{\textbf{Moralité :}} \alert{aucune propriété
  non-triviale} de la fonction $\varphi^{(1)}_e$ calculée par un
programme \alert{n'est décidable} en regardant le programme $e$.

\bigskip

Exemples :
\begin{itemize}
\item $\{e \in \mathbb{N} : \varphi^{(1)}_e(0){\downarrow}\}$ n'est
  pas décidable ($\Rightarrow$ Rice \alert{généralise}
  l'indécidabilité du pb. de l'arrêt).
\item $\{e \in \mathbb{N} : \varphi^{(1)}_e \text{~totale}\}$ n'est
  pas décidable.
\item $\{e \in \mathbb{N} : \forall
  n.\,(\varphi^{(1)}_e(n){\downarrow} \,\Rightarrow\,
  \varphi^{(1)}_e(n) = 0)\}$ n'est pas décidable.
\end{itemize}

\end{frame}
%
\begin{frame}
\frametitle{Le théorème de Rice : preuve par théorème de récursion}

{\footnotesize $\textbf{R}^{(1)} = \{f \colon
  \mathbb{N}\dasharrow\mathbb{N} : f\text{~calculable}\}$\par}

\smallskip

\itempoint\textbf{Théorème} (Rice) : si $F \subseteq
\textbf{R}^{(1)}$ est un ensemble de fonctions partielles tel que
$\Phi^{-1}(F) := \{e \in \mathbb{N} : \varphi^{(1)}_e \in F\}$ est
\emph{décidable}, alors $F = \varnothing$ ou $F = \textbf{R}^{(1)}$.

\bigskip

{\footnotesize La preuve est très analogue à celle de l'indécidabilité
  du problème de l'arrêt.\par}

\smallskip

\underline{Preuve :} Supposons par l'absurde $\Phi^{-1}(F)$ décidable
avec $F \neq \varnothing$ et $F \neq \textbf{R}^{(1)}$.  Soient $f \in
F$ et $g \not\in F$.  Soit
\[
h(e,x) := \left\{
\begin{array}{ll}
f(x)&\text{~si~}e\not\in \Phi^{-1}(F)\\
g(x)&\text{~si~}e\in \Phi^{-1}(F)\\
\end{array}
\right.
\]
Alors $h \colon \mathbb{N}^2 \dasharrow \mathbb{N}$ est calculable par
hypothèse (on peut décider si $e\in \Phi^{-1}(F)$).  Par le théorème
de récursion de Kleene (transp. \ref{kleene-recursion-theorem}), il
existe $e$ tel que
\[\varphi^{(1)}_e(x) = h(e,x)\]

Si $e \in \Phi^{-1}(F)$ alors $h(e,x) = g(x)$ pour tout $x$, donc
$\Phi(e) = g$ donc $e \not\in \Phi^{-1}(F)$, une contradiction.  Si $e
\not\in \Phi^{-1}(F)$ alors $h(e,x) = f(x)$ pour tout $x$, donc
$\Phi(e) = f$ donc $e \in \Phi^{-1}(F)$, une contradiction.\qed

\end{frame}
%
\begin{frame}
\frametitle{Réductions : introduction}

\itempoint Situation typique : on veut montrer qu'une question $D$
(« problème de décision », souvent déjà semi-décidable) est
indécidable.  Ceci se fait typiquement en \alert{réduisant le problème
  de l'arrêt} à $D$, c'est-à-dire :

\bigskip

\textcolor{teal}{« Supposons par l'absurde que $D$ soit décidable,
  c'est-à-dire que j'ai un algorithme qui répond à la question $D$
  (comprendre : “$n\in D$ ?”).}

\smallskip

\textcolor{teal}{Je montre qu'\alert{en utilisant cet algorithme} je
  peux résoudre le problème de l'arrêt.}

\smallskip

\textcolor{teal}{Ceci est une contradiction (car le problème de
  l'arrêt est indécidable),}

\textcolor{teal}{donc $D$ est indécidable. »}

\bigskip

\itempoint Les notions de réduction formalisent cet argument :
intuitivement,

\centerline{« $A$ se réduit à $B$ »}

\centerline{signifie}

\centerline{« si $B$ est décidable alors $A$ est décidable »}

\centerline{(mais constructivement)}

\end{frame}
%
\begin{frame}
\frametitle{Le théorème de Rice : preuve par réduction (1/2)}

{\footnotesize $\textbf{R}^{(1)} = \{f \colon
  \mathbb{N}\dasharrow\mathbb{N} : f\text{~calculable}\}$\par}

\smallskip

\itempoint\textbf{Théorème} (Rice) : si $F \subseteq
\textbf{R}^{(1)}$ est tel que $F \neq \varnothing$ et $F \neq
\textbf{R}^{(1)}$, alors $\Phi^{-1}(F) := \{e \in \mathbb{N} :
\varphi^{(1)}_e \in F\}$ \emph{n'est pas décidable}.

\bigskip

\underline{Preuve :} Soit $F \subseteq \textbf{R}^{(1)}$ avec $F \neq
\varnothing$ et $F \neq \textbf{R}^{(1)}$.  Quitte à remplacer $F$
par $\complement F$, o.p.s. ${\uparrow} \not\in F$ où $\uparrow$ est
la fonction nulle part définie.  Soit $f\in F$ où $f =
\varphi^{(1)}_a$.

\smallskip

Pour $(e,x) \in \mathbb{N}^2$, considérons l'algorithme suivant,
prenant en entrée $m \in \mathbb{N}$ :
\begin{itemize}
\item simuler $\varphi^{(1)}_e(x)$ avec la machine universelle, puis,
  si l'exécution termine,
\item calculer $f(m) = \varphi^{(1)}_a(m)$ et (si l'exécution termine)
  renvoyer sa valeur.
\end{itemize}

\smallskip

Soit $b(e,x)$ le code de l'algorithme qu'on vient de décrire :

\centerline{$\varphi^{(1)}_{b(e,x)} = f$ si $\varphi^{(1)}_e(x)\downarrow$
\quad et\quad
$\varphi^{(1)}_{b(e,x)} = {\uparrow}$ si $\varphi^{(1)}_e(x)\uparrow$}

notamment $\varphi^{(1)}_{b(e,x)} \in F$ ssi $\varphi^{(1)}_e(x)\downarrow$
\textcolor{brown}{($\leftarrow$ c'est là la réduction)}.\hfill  …/…

\end{frame}
%
\begin{frame}
\frametitle{Le théorème de Rice : preuve par réduction (2/2)}

{\footnotesize $\textbf{R}^{(1)} = \{f \colon
  \mathbb{N}\dasharrow\mathbb{N} : f\text{~calculable}\}$ ; on a
  supposé $F \subseteq \textbf{R}^{(1)}$ avec ${\uparrow}\not\in F$
  et $f \in F$\par}

\smallskip

On a construit (transp. précédent) un $b(e,x)$, avec $b \colon
\mathbb{N}^2 \to \mathbb{N}$ calculable (même p.r.) tel que
$\varphi^{(1)}_{b(e,x)} \in F$ ssi $\varphi^{(1)}_e(x)\downarrow$,
c'est-à-dire
\[
b(e,x) \in \Phi^{-1}(F) \;\Longleftrightarrow\; (e,x) \in \mathscr{H}
\]
où $\mathscr{H} := \{(e,x)\in\mathbb{N}^2 : \varphi_e(x)\downarrow\}$
est le problème de l'arrêt.

\medskip

Si $\Phi^{-1}(F)$ était décidable, alors $\mathscr{H}$ le serait aussi, par
l'algorithme :
\begin{itemize}
\item donnés $e,x$, calculer $b(e,x)$, décider si $b(e,x) \in \Phi^{-1}(F)$,
\item si oui, répondre « oui », sinon répondre « non ».
\end{itemize}

\smallskip

Or $\mathscr{H}$ n'est pas décidable, donc $\Phi^{-1}(F)$ non plus.\qed

\bigskip

On dit qu'on a \alert{réduit le problème de l'arrêt} à $\Phi^{-1}(F)$
(\alert{via} la fonction $b$).

\end{frame}
%
\begin{frame}
\frametitle{Réduction « many-to-one »}

\textbf{Définition :} Si $A,B\subseteq\mathbb{N}$, on note $A
\mathrel{\leq_\mathrm{m}} B$ lorsqu'il existe $\rho \colon \mathbb{N}
\to \mathbb{N}$ \emph{calculable totale} telle que
\[
\rho(m) \in B \;\Longleftrightarrow\; m \in A
\]
{\footnotesize (c'est-à-dire $A = \rho^{-1}(B)$)}.

\bigskip

\textcolor{blue}{\textbf{Intuitivement :}} si j'ai un gadget qui
répond à la question “$n \in B$ ?”, je peux répondre à la question “$m
\in A$ ?” en transformant $m$ en $\rho(m) =: n$ et en utilisant le gadget
{\footnotesize (une seule fois, à la fin)}.

\bigskip

\textbf{Clairement}, si $A \mathrel{\leq_\mathrm{m}} B$ avec $B$
décidable (resp. semi-décidable), alors $A$ est décidable
(resp. semi-décidable).

\smallskip

\emph{Notamment}, si $\mathscr{H} \mathrel{\leq_\mathrm{m}} D$ alors
$D$ \emph{n'est pas} décidable.

\bigskip

{\footnotesize La relation $\mathrel{\leq_\mathrm{m}}$ est réflexive
  et transitive (c'est un « préordre ») ; la relation
  $\mathrel{\equiv_\mathrm{m}}$ définie par $A
  \mathrel{\equiv_\mathrm{m}} B$ ssi $A \mathrel{\leq_\mathrm{m}} B$
  et $B \mathrel{\leq_\mathrm{m}} A$ est une relation d'équivalence,
  les classes pour laquelle s'appellent « degrés many-to-one » et sont
  partiellement ordonnés par $\mathrel{\leq_\mathrm{m}}$.\par}

\end{frame}
%
\begin{frame}
\frametitle{Réduction de Turing : présentation informelle}

\textbf{Informellement :} Si $A,B\subseteq\mathbb{N}$, on note $A
\mathrel{\leq_\mathrm{T}} B$ s'il existe un algorithme qui
\begin{itemize}
\item prend en entrée $m \in \mathbb{N}$,
\item peut à tout moment demander à savoir si $n \in B$
  (\textcolor{teal}{« interroger l'oracle »}),
\item termine en temps fini,
\item et renvoie « oui » si $m \in A$, et « non » si $m \not\in A$.
\end{itemize}

\bigskip

\textcolor{blue}{\textbf{Intuitivement :}} à la différence de la
réduction many-to-one où on ne peut poser la question “$n \in B$ ?”
que sur une seule valeur $\rho(n)$ à la fin du calcul, ici on peut
interroger l'oracle de façon libre et illimitée (mais finie !) au
cours de l'algorithme.

\bigskip

La relation $A \mathrel{\leq_\mathrm{T}} B$ est beaucoup plus faible
que $A \mathrel{\leq_\mathrm{m}} B$.

\smallskip

{\footnotesize Par exemple, $(\complement B) \mathrel{\leq_\mathrm{T}}
  B$ pour tout $B\subseteq\mathbb{N}$ (savoir décider “$n \in B$ ?”
  permet évidemment de décider “$n \not\in B$ ?”), alors que
  $(\complement \mathscr{H}) \mathrel{\not\leq_\mathrm{m}}
  \mathscr{H}$ car $\complement \mathscr{H}$ n'est pas
  semi-décidable.\par}

\bigskip

\textcolor{brown}{Mais comment formaliser cette « interrogation » ?}

\end{frame}
%
\begin{frame}
\frametitle{Réduction de Turing : formalisation(s) possible(s)}

Comment définir $A \mathrel{\leq_\mathrm{T}} B$ pour $A, B \subseteq
\mathbb{N}$ ?  {\footnotesize (I.e., « $A$ est calculable en
  utilisant $B$ ».)}

\bigskip

\textbf{Formalisation 1 :} la fonction indicatrice $\mathbf{1}_A$
de $A$ appartient à la plus petite classe de fonctions qui contient
les projections, les constantes, la fonction successeur \alert{et la
  fonction indicatrice $\mathbf{1}_B$ de $B$} et stable par
composition, récursion primitive et opérateur $\mu$.

\bigskip

\textbf{Formalisation 2 :} il existe une fonction calculable qui prend
en entrée $m \in \mathbb{N}$ et une liste $\dbllangle \langle n_0,
\mathbf{1}_B(n_0)\rangle, \ldots, \langle n_k,
\mathbf{1}_B(n_k)\rangle \dblrangle$ de réponses à des questions “$n
\in B$ ?”, et qui (si ces réponses cont correctes !) termine et renvoie
\begin{itemize}
\item soit une réponse finale à la question “$m \in A$ ?” (disons
  encodée comme $\langle 0, \mathbf{1}_A(m)\rangle$),
\item soit une nouvelle interrogation “$n \in B$ ?” (disons encodée
  comme $\langle 1, n\rangle$),
\end{itemize}
de sorte que si on commence par $k=0$ et qu'on ajoute à chaque fois la
réponse correcte $\langle n_{k+1}, \mathbf{1}_B(n_{k+1})\rangle$ à
l'interrogation $\langle 1, n_{k+1}\rangle$ posée, alors la fonction
finit par produite la réponse finale correcte ($\langle 0,
\mathbf{1}_A(m)\rangle$).

\end{frame}
%
\begin{frame}
\frametitle{Réduction de Turing : quelques propriétés}

\textbf{Clairement}, si $A \mathrel{\leq_\mathrm{T}} B$ avec $B$
décidable, alors $A$ est décidable.

{\footnotesize (Ceci \alert{ne vaut pas} pour
  « semi-décidable ».)\par}

\smallskip

\emph{Notamment}, si $\mathscr{H} \mathrel{\leq_\mathrm{T}} D$ alors
$D$ \emph{n'est pas} décidable.

\bigskip

{\footnotesize

La relation $\mathrel{\leq_\mathrm{T}}$ est réflexive et transitive
(c'est un « préordre ») ; la relation $\mathrel{\equiv_\mathrm{T}}$
définie par $A \mathrel{\equiv_\mathrm{T}} B$ ssi $A
\mathrel{\leq_\mathrm{T}} B$ et $B \mathrel{\leq_\mathrm{T}} A$ est
une relation d'équivalence, les classes pour laquelle s'appellent
« degrés de Turing » et sont partiellement ordonnés par
$\mathrel{\leq_\mathrm{T}}$.

\medskip

Comme $A \mathrel{\leq_\mathrm{m}} B$ implique $A
\mathrel{\leq_\mathrm{T}} B$, chaque degré de Turing est une réunion
de degrés many-to-one (la relation d'équivalence
$\mathrel{\equiv_\mathrm{T}}$ est plus grossière que
$\mathrel{\equiv_\mathrm{m}}$).

\medskip

Les parties décidables de $\mathbb{N}$ forment le plus petit degré de
Turing, souvent noté $\mathbf{0}$.  Le degré de Turing de
$\mathscr{H}$ est noté $\mathbf{0'}$.  (Il existe des ensembles de
degré strictement compris entre $\mathbf{0}$ et $\mathbf{0'}$, même
des ensembles semi-décidables, mais il semble qu'aucun n'apparaît
naturellement.)

\par}

\end{frame}
%
\section{Le \texorpdfstring{$\lambda$}{lambda}-calcul non typé}
\begin{frame}
\frametitle{Le $\lambda$-calcul : aperçu}

Le \textbf{$\lambda$-calcul non typé} manipule des expressions du type
\[
\begin{array}{c}
\lambda x.\lambda y.\lambda z.((xz)(yz))\\
\lambda f.\lambda x.f(f(f(f(f(fx)))))\\
(\lambda x.(xx))(\lambda x.(xx))\\
\end{array}
\]

\bigskip

Ces expressions s'appelleront des \textbf{termes} du $\lambda$-calcul.

\bigskip

Il faut comprendre intuitivement qu'un terme représente une sorte de
fonction qui prend une autre telle fonction en entrée et renvoie une
autre telle fonction.

\bigskip

Deux constructions fondamentales :
\begin{itemize}
\item\textbf{application} : $(PQ)$ : appliquer la fonction $P$ à la
  fonction $Q$ ;
\item\textbf{abstraction} : $\lambda v.E$ : créer la fonction qui
  prend un argument et le remplace pour $v$ dans l'expression $E$
  (\textcolor{teal}{en gros} $v \mapsto E$).
\end{itemize}

\end{frame}
%
\begin{frame}
\frametitle{Le $\lambda$-calcul : termes}

\itempoint Un \textbf{terme} du $\lambda$-calcul est (inductivement) :
\begin{itemize}
\item une \textbf{variable} ($a$, $b$, $c$... en nombre illimité),
\item une \textbf{application} $(PQ)$ où $P$ et $Q$ sont deux termes,
\item une \textbf{abstraction} $\lambda v.E$ où $v$ est une variable
  et $E$ un terme ; on dira que la variable $v$ est \textbf{liée} dans
  $E$ par ce $\lambda$.
\end{itemize}

\medskip

\itempoint Conventions d'écriture :
\begin{itemize}
\item l'application \alert{n'est pas associative} : on parenthèse
  implicitement vers la gauche : « $xyz$ » signifie « $((xy)z)$ » ;
\item abréviation de plusieurs $\lambda$ : on note « $\lambda uv.E$ »
  pour « $\lambda u. \lambda v. E$ » ;
\item l'abstraction est moins prioritaire que l'application :
  « $\lambda x.xy$ » signifie $\lambda x.(xy)$ \alert{pas} $(\lambda
  x.x)y$.
\end{itemize}

\medskip

\itempoint Une variable non liée est dite \textbf{libre} : $(\lambda
  x.x)\textcolor{red}{x}$ (le dernier $\textcolor{red}{x}$ est libre).

\itempoint Un terme sans variable libre est dit \textbf{clos}.

\itempoint Les variables liées sont muettes : $\lambda x.x \equiv \lambda
  y.y$, comprendre $\mathord{\tikz[remember picture, baseline =
      (binder.base), inner sep = 0pt] {\node (binder)
      {\strut$\lambda\bullet$};}}.\mathord{\tikz[remember picture, baseline =
      (bindee.base), inner sep = 0pt] {\node (bindee) {\strut$\bullet$};}}$.
\begin{tikzpicture}[remember picture, overlay]
\draw [->, >=stealth, thick] (bindee.north) -- ($(bindee.north)+(0pt,8pt)$) -- ($(binder.north)+(0pt,8pt)$) -- (binder.north);
\end{tikzpicture}

\end{frame}
%
\begin{frame}
\frametitle{Le $\lambda$-calcul : variables liées}

On appelle \textbf{$\alpha$-conversion} le renommage des variables
liées : ces termes sont considérés comme équivalents.
\begin{itemize}
\item $\lambda x.x \equiv \lambda y.y$ et $\lambda xyz.((xz)(yz))
  \equiv \lambda uvw.((uw)(vw))$
\item Attention à \alert{ne pas capturer} de variable libre : $\lambda
  y.xy \mathrel{\textcolor{red}{\not\equiv}} \lambda x.xx$.
\item En cas de synonymie, la variable est liée par le $\lambda$ le
  \alert{plus intérieur} pour ce nom ($\cong$ portée lexicale) :
  $\lambda x. \lambda x. x \equiv \lambda x. \lambda v. v
  \mathrel{\textcolor{red}{\not\equiv}} \lambda u. \lambda x. u$.
\item Mieux vaut ne pas penser aux termes typographiquement, mais à
  chaque variable liée comme un \emph{pointeur vers la
  $\lambda$-abstraction qui la lie} :
\[
\lambda x. (\lambda y. y (\lambda z. z)) (\lambda z. x z)
\equiv
\textcolor{red}{\lambda\bullet}. (\textcolor{yellow}{\lambda\bullet}. \textcolor{yellow}{\bullet} (\textcolor{green}{\lambda\bullet}. \textcolor{green}{\bullet})) (\textcolor{blue}{\lambda\bullet}. \textcolor{red}{\bullet} \textcolor{blue}{\bullet})
\]
\item Autre convention possible : \textbf{indices de De Bruijn} :
  remplacer les variables liées par le numéro du $\lambda$ qui la lie,
  en comptant du plus intérieur ($1$) vers le plus extérieur :
\[
\lambda x. (\lambda y. y (\lambda z. z)) (\lambda z. x z)
\equiv
\lambda. (\lambda. 1 (\lambda. 1)) (\lambda. 2 1)
\]
deux termes sont $\alpha$-équivalents ssi leur écriture avec indice de
De Bruijn est identique.
\end{itemize}

\end{frame}
%
\begin{frame}
\frametitle{Le $\lambda$-calcul : $\beta$-réduction}

{\footnotesize On travaille désormais sur des termes à
  $\alpha$-équivalence près.\par}

\bigskip

\itempoint Un \textbf{redex} (« reducible expression ») est un terme
de la forme $(\lambda v. E)T$.

Son \textbf{réduit} est le terme $E[v\backslash T]$ obtenu par
remplacement de $T$ pour $v$ dans $E$, en évitant tout conflit de
variables.

\medskip

Exemples :
\begin{itemize}
\item $(\lambda x.xx)y \rightarrow yy$
\item $(\lambda x.xx)(\lambda x.xx) \rightarrow (\lambda x.xx)(\lambda
  x.xx)$ (est son propre réduit)
\item $(\lambda xy.x)z \rightarrow \lambda y.z$ (car $\lambda xy.x$
  abrège $\lambda x.\lambda y.x$)
\item $(\lambda xy.x)y \rightarrow \lambda y_1.y$ (attention au
  conflit de variable !)
\item $(\lambda x.\lambda x.x)y \rightarrow \lambda x.x$ (car $\lambda
  x.\lambda x.x \equiv \lambda z.\lambda x.x$ : le $\lambda$ extérieur
  ne lie rien)
\end{itemize}

\bigskip

\itempoint Un terme n'ayant \alert{pas de redex en sous-expression}
est dit en \textbf{forme ($\beta$-)normale}.\quad  Ex. : $\lambda
xyz.((xz)(yz))$.

\smallskip

\itempoint On appelle \textbf{$\beta$-réduction} le remplacement en
sous-expression d'un \textcolor{purple}{redex} par son
\textcolor{olive}{réduit}.\quad Ex. : $\lambda
x. \textcolor{purple}{(\lambda y. y (\lambda z. z)) (\lambda z. x z)}
\rightarrow \lambda x. \textcolor{olive}{(\lambda z. x z)(\lambda
  z. z)}$.

\end{frame}
%
\begin{frame}
\frametitle{Le $\lambda$-calcul : normalisation par $\beta$-réductions}

On note :
\begin{itemize}
\item $T \rightarrow T'$ (ou $T \rightarrow_\beta T'$) si $T'$
  s'obtient par $\beta$-réduction d'un redex de $T$.
\item $T \twoheadrightarrow T'$ (ou $T \twoheadrightarrow_\beta T'$)
  si $T'$ s'obtient par une suite finie de $\beta$-réductions ($T =
  T_0 \rightarrow \cdots \rightarrow T_n = T'$, y compris $n=0$ soit
  $T'=T$).
\item $T$ est \textbf{faiblement normalisable} lorsque $T
  \twoheadrightarrow T'$ avec $T'$ en forme normale (\alert{une
    certaine} suite de $\beta$-réductions termine).
\item $T$ est \textbf{fortement normalisable} lorsque \alert{toute}
  suite de $\beta$-réductions termine (sur un terme en forme normale).
\end{itemize}

\bigskip

Exemples :
\begin{itemize}
\item $(\lambda x.xx)(\lambda x.xx)$ n'est pas faiblement
  normalisable (la $\beta$-réduction boucle).
\item $(\lambda uz.z)((\lambda x.xx)(\lambda x.xx))$ n'est
  pas fortement normalisable mais il est faiblement normalisable
  $\rightarrow \lambda z.z$.
\item $(\lambda uz.u)((\lambda t.t)(\lambda x.xx))$ est fortement
  normalisable $\twoheadrightarrow \lambda zx.xx$.
\end{itemize}

\end{frame}
%
\begin{frame}
\frametitle{Le $\lambda$-calcul : confluence et choix d'un redex}

\itempoint\textbf{Théorème} (Church-Rosser) : si $T \twoheadrightarrow
T'_1$ et $T \twoheadrightarrow T'_2$ alors il existe $T''$ tel que
$T'_1 \twoheadrightarrow T''$ et $T'_2 \twoheadrightarrow T''$.

\smallskip

En particulier, si $T'_1, T'_2$ sont en forme normale, alors $T'_1
\equiv T'_2$ (unicité de la normalisation).

\bigskip

Pour \alert{éviter} ce théorème, on va faire un choix simple de redex
à réduire :

\itempoint On appelle \textbf{redex extérieur gauche} d'un
$\lambda$-terme le redex dont le $\lambda$ est \alert{le plus à
  gauche}.  Exemples : $\lambda x.x(\textcolor{purple}{(\lambda
  y.y)x})$ ; $\lambda x.\textcolor{purple}{(\lambda y.(\lambda
  z.z)y)x}$.

\medskip

\itempoint On écrira $T \rightarrow_{\mathsf{lft}} T'$ lorsque $T'$
s'obtient par $\beta$-réduction du redex extérieur gauche, et $T
\twoheadrightarrow_{\mathsf{lft}} T'$ pour une suite de telles
réductions.

\bigskip

On peut montrer (mais on évitera d'utiliser) :

\itempoint\textbf{Théorème} (Curry \&al) : si $T \twoheadrightarrow
T'$ avec $T'$ en forme normale, alors $T
\twoheadrightarrow_{\mathsf{lft}} T'$ (i.e., la réduction ext.
gauche normalise les termes faiblement normalisables).

\end{frame}
%
\begin{frame}
\frametitle{Réduction extérieure gauche : exemples}

{\footnotesize Divers noms utilisés : « réduction en ordre normal »,
  « réduction gauche », etc.\par}

\bigskip

On a noté $T \twoheadrightarrow_{\mathsf{lft}} T'$ lorsque $T'$
s'obtient par une succession de $\beta$-réductions à chaque fois du
redex dont le $\lambda$ est le plus à gauche.

\bigskip

Exemples :
\begin{itemize}
\item $\textcolor{purple}{(\lambda x.xx)(\lambda x.xx)}
  \rightarrow_{\mathsf{lft}} (\lambda x.xx)(\lambda x.xx)
  \rightarrow_{\mathsf{lft}} \cdots$ (boucle)
\item $(\lambda uz.z)((\lambda x.xx)(\lambda x.xx)) =
  \textcolor{purple}{(\lambda u.\lambda z.z)((\lambda x.xx)(\lambda
    x.xx))} \rightarrow_{\mathsf{lft}} \lambda z.z$
\item $(\lambda uz.u)((\lambda t.t)(\lambda x.xx)) =
  \textcolor{purple}{(\lambda u.\lambda z.u)((\lambda t.t)(\lambda
    x.xx))} \rightarrow_{\mathsf{lft}} \lambda
  z.(\textcolor{purple}{(\lambda t.t)(\lambda x.xx)})
  \rightarrow_{\mathsf{lft}} \lambda z.\lambda x.xx = \lambda zx.xx$
\end{itemize}

\bigskip

Intérêt :
\begin{itemize}
\item cette stratégie de réduction est \alert{déterministe},
\item (Curry \&al :) si (« terme faiblement normalisant ») une
  réduction quelconque termine sur une forme normale, alors
  $\twoheadrightarrow_{\mathsf{lft}}$ le fait.
\end{itemize}

\end{frame}
%
\begin{frame}
\frametitle{Simulation du $\lambda$-calcul par les fonctions récursives}

\itempoint On peut coder un terme du $\lambda$-calcul sous forme
d'entiers naturels.

\bigskip

\itempoint La fonction $T \mapsto 1$ qui à un terme $T$ associe $0$ si
$T$ est en forme normale et $1$ si non, \textbf{est p.r.}

\medskip

\itempoint La fonction $T \mapsto T'$ qui à un terme $T$ associe sa
réduction extérieure gauche \textbf{est p.r.}

\medskip

\itempoint Conséquence : la fonction $(n,T) \mapsto T^{(n)}$ qui à
$n\in\mathbb{N}$ et un terme $T$ associe le terme obtenu après $n$
réductions extérieures gauches \textbf{est p.r.}

\medskip

\itempoint La fonction qui à $T$ associe la forme normale (et/ou le
nombre d'étapes d'exécution) \alert{si la réduction extérieure gauche
  termine}, et $\uparrow$ (non définie) si elle ne termine pas, est
\textbf{générale récursive}.

\bigskip

\textcolor{blue}{\textbf{Moralité :}} les fonctions récursives peuvent
simuler la réduction extérieure gauche du $\lambda$-calcul
{\footnotesize (ou n'importe quelle autre réduction, mais on se
  concentre sur celle-ci)}.

\end{frame}
%
\begin{frame}
\frametitle{Entiers de Church}

On définit les termes en forme normale $\overline{n} := \lambda
fx.f^{\circ n}(x)$ pour $n\in\mathbb{N}$, c-à-d :
\begin{itemize}
\item $\overline{0} := \lambda fx.x$
\item $\overline{1} := \lambda fx.fx$
\item $\overline{2} := \lambda fx.f(fx)$
\item $\overline{3} := \lambda fx.f(f(fx))$, etc.
\end{itemize}

{\footnotesize Intuitivement, $\overline{n}$ prend une fonction
  $f$ et renvoie sa $n$-ième itérée.\par}

\medskip

\itempoint Posons $A := \lambda mfx.f(mfx) = \lambda m.\lambda
f.\lambda x.f(mfx)$

Alors
\[
\begin{aligned}
A\overline{n} &= (\lambda m.\lambda f.\lambda x.f(mfx))(\lambda
g.\lambda y.g^{\circ n}(y))\\
& \rightarrow_{\mathsf{lft}} \lambda f.\lambda
x.f(((\lambda g.\lambda y.g^{\circ n}(y)))fx)\\
&\rightarrow_{\mathsf{lft}}\rightarrow_{\mathsf{lft}} \lambda f.\lambda
x.f(f^{\circ n}(x)) = \lambda f.\lambda x.f^{\circ(n+1)}(x) = \overline{n+1}
\end{aligned}
\]

\end{frame}
%
\begin{frame}
\frametitle{Calculs dans le $\lambda$-calcul : une convention}

On dira qu'une fonction $f\colon \mathbb{N}^k \dasharrow \mathbb{N}$
est \textbf{représentable par un $\lambda$-terme} lorsqu'il existe un
terme clos $t$ tel que, pour tous $x_1,\ldots,x_k \in \mathbb{N}$ :
\begin{itemize}
\item si $f(x_1,\ldots,x_k){\downarrow}=y$ alors
  $t\overline{x_1}\cdots\overline{x_k}
  \twoheadrightarrow_{\mathsf{lft}} \overline{y}$,
\item si $f(x_1,\ldots,x_k){\uparrow}$ alors
  $t\overline{x_1}\cdots\overline{x_k} \rightarrow_{\mathsf{lft}}
  \cdots$ ne termine pas,
\end{itemize}
où $\overline{z}$ désigne l'entier de Church associé
à $z\in\mathbb{N}$.

\bigskip

Exemples :
\begin{itemize}
\item $\lambda mfx.f(mfx)$ représente $m \mapsto m+1$
  (transp. précédent),
\item $\lambda mnfx.nf(mfx)$ représente $(m,n) \mapsto m+n$,
\item $\lambda mnf.n(mf)$ représente $(m,n) \mapsto mn$ {\footnotesize
  (itérer $n$ fois l'itérée $m$-ième)},
\item $\lambda mn.nm$ représente $(m,n) \mapsto m^n$ {\footnotesize
  (itérer $n$ fois l'itération $m$-ième)}.
\item $\lambda mnp.p(\lambda y.n)m$ représente $(m,n,p) \mapsto
  \left\{\begin{array}{ll}m&\text{~si~}p=0\\n&\text{~si~}p\geq
  1\end{array}\right.$\\{\footnotesize (itérer $p$ fois « remplacer
    par $n$ »)}.
\end{itemize}

\end{frame}
%
\begin{frame}
\frametitle{Représentation des fonctions p.r. : cas faciles}

{\footnotesize (Cf. transp. \ref{primitive-recursive-definition}.)\par}

Fonction p.r. facilement représentables par un $\lambda$-terme :
\begin{itemize}
\item $\lambda x_1\cdots x_k.x_i$ représente $(x_1,\ldots,x_k) \mapsto x_i$ ;
\item $\lambda x_1\cdots x_k.\overline{c}$ représente
  $(x_1,\ldots,x_k) \mapsto c$ ;
\item $A := \lambda mfx.f(mfx)$ représente $x \mapsto x+1$ ;
\item si $v_1,\ldots,v_\ell$ représentent $g_1,\ldots,g_\ell$ et $w$
  représente $h$, alors $\lambda x_1\cdots x_k.w(v_1 x_1\cdots
  x_k)\cdots (v_\ell x_1\cdots x_k)$ représente $(x_1,\ldots,x_k)
  \mapsto h(g_1(x_1,\ldots,x_k),\ldots, g_\ell(x_1,\ldots,x_k))$ ;
\item si $v$ représente $g$ et $w$ représente $h$, alors
\[
\lambda x_1\cdots x_k z.z(wx_1\cdots x_k)(vx_1\cdots x_k)
\]
représente $f$ définie par la récursion primitive
\[
\begin{aligned}
f(x_1,\ldots,x_k,0) &= g(x_1,\ldots,x_k)\\
f(x_1,\ldots,x_k,z+1) &= h(x_1,\ldots,x_k,f(x_1,\ldots,x_k,z))
\end{aligned}
\]
\alert{mais} on veut $f(x_1,\ldots,x_k,z+1) =
h(x_1,\ldots,x_k,f(x_1,\ldots,x_k,z),\alert{z})$...?
\end{itemize}

\end{frame}
%
\begin{frame}
\frametitle{Représentation des couples d'entiers}

{\footnotesize (Oublions $x_1,\ldots,x_k$ pour ne pas alourdir les
  notations.)\par}

Comment passer de
\[
\left\{
\begin{aligned}
f(0) &= g\\
f(z+1) &= h(f(z))
\end{aligned}
\right.
\quad\text{~à~}\quad
\left\{
\begin{aligned}
f(0) &= g\\
f(z+1) &= h(f(z),z)
\end{aligned}
\right.
\quad\text{~?}
\]
On voudrait définir
\[
\tilde f(z) = (f(z),z)
\quad\text{~soit~}\quad
\left\{
\begin{aligned}
\tilde f(0) &= (g,0)\\
\tilde f(z+1) &= \tilde h(\tilde f(z))
\end{aligned}
\right.
\;\text{~où~}\;
\tilde h(y,z) = (h(y,z), z+1)
\]

\bigskip

On va définir (temporairement ?)
\[
\begin{aligned}
\overline{m,n} &:= \lambda fgx.f^{\circ m}(g^{\circ n}(x))
\quad\text{~si~}m,n\in\mathbb{N}\\
\Pi &:= \lambda mnfgx.(mf)(ngx)
\quad\text{~donc~}\Pi\overline{m}\,\overline{n} \twoheadrightarrow_{\mathsf{lft}} \overline{m,n}\\
\pi_1 &:= \lambda pfx.pf(\lambda z.z)x
\quad\text{~donc~}\pi_1\overline{m,n} \twoheadrightarrow_{\mathsf{lft}} \overline{m}\\
\pi_2 &:= \lambda pgx.p(\lambda z.z)gx
\quad\text{~donc~}\pi_2\overline{m,n} \twoheadrightarrow_{\mathsf{lft}} \overline{n}\\
\end{aligned}
\]

\end{frame}
%
\begin{frame}
\frametitle{Représentation de la récursion primitive générale}

Maintenant qu'on a une représentation des couples d'entiers naturels
dans le $\lambda$-calcul donnée par $\Pi$ (formation de paires) et
$\pi_1,\pi_2$ (projections).

\bigskip

\itempoint Si $v$ représente $g\colon \mathbb{N}^k \dasharrow
\mathbb{N}$ et $w$ représente $h\colon \mathbb{N}^{k+2} \dasharrow
\mathbb{N}$, alors $f\colon \mathbb{N}^{k+1} \dasharrow \mathbb{N}$
est représentée par
\[
\lambda x_1\cdots x_k z.
\pi_1(z(\lambda p.\Pi(w x_1\cdots x_k (\pi_1 p)(\pi_2 p))A(\pi_2 p))(\Pi(vx_1\cdots x_k)\overline{0}))
\]
où
\[
\begin{aligned}
f(x_1,\ldots,x_k,0) &= g(x_1,\ldots,x_k)\\
f(x_1,\ldots,x_k,z+1) &= h(x_1,\ldots,x_k,f(x_1,\ldots,x_k,z),z)
\end{aligned}
\]
(toujours avec $A := \lambda mfx.f(mfx)$).

\bigskip

{\footnotesize D'autres encodages des paires sont possibles et
  possiblement plus simples, p.ex., $\Pi := \lambda rsa.ars$ et $\pi_1
  := \lambda p.p(\lambda rs.r)$ et $\pi_2 := \lambda p.p(\lambda
  rs.s)$ (fonctionnent sur plus que les entiers de Church).\par}

\bigskip

Bref, (au moins) \alert{les fonctions p.r. sont représentables par
  $\lambda$-termes}.

\end{frame}
%
\begin{frame}
\frametitle{Le combinateur $\mathsf{Y}$ de Curry}

\itempoint Pour représenter toutes les fonctions récursives, on va
implémenter les appels récursifs dans le $\lambda$-calcul.

\bigskip

\itempoint Pour ça, on va utiliser la même idée que le théorème de
récursion de Kleene
(transp. \ref{kleene-recursion-theorem-p-r-version}).

\bigskip

Posons
\[
\mathsf{Y} := \lambda f. ((\lambda x.f(x x)) (\lambda x.f(x x)))
\]

Idée :
\[
\begin{aligned}
\mathsf{Y} &:= \lambda f. ((\lambda x.f(x x)) (\lambda x.f(x x)))\\
&\rightarrow \lambda f. f((\lambda x.f(x x)) (\lambda x.f(x x)))\\
&\rightarrow \lambda f. f(f((\lambda x.f(x x)) (\lambda x.f(x x))))
\rightarrow \cdots
\end{aligned}
\]

\itempoint Le terme (non normalisable !) $\mathsf{Y}$
“\textbf{recherche}” un point fixe de son argument.

\bigskip

\itempoint Permet d'implémenter les appels récursifs (transp. suivant).

\end{frame}
%
\begin{frame}
\frametitle{Récursion avec le combinateur $\mathsf{Y}$}

\itempoint On veut implémenter une définition par appels récursifs
dans le $\lambda$-calcul, $\texttt{let~rec~}h = (\ldots h\ldots)$,
disons $\texttt{let~rec~}h = E$ où $E = (\ldots h\ldots)$ est un terme
faisant intervenir $h$.

\smallskip

{\footnotesize L'idée est comme dans le
  transp. \ref{recursion-from-kleene-recursion-theorem}.\par}

\medskip

\itempoint On considère le terme :
\[
\begin{aligned}
\mathsf{Y}(\lambda h.E)
&= (\lambda f. ((\lambda x.f(x x)) (\lambda x.f(x x))))(\lambda h.E)\\
&\rightarrow (\lambda x.(\lambda h.E)(x x)) (\lambda x.(\lambda h.E)(x x))
=: h_{\mathrm{fix}}\\
&\rightarrow (\lambda h.E)((\lambda x.(\lambda h.E)(x x)) (\lambda x.(\lambda h.E)(x x)))
= (\lambda h.E)(h_{\mathrm{fix}})\\
&\rightarrow E[h\backslash h_{\mathrm{fix}}]
\text{~(substitution de $h_{\mathrm{fix}}$ pour $h$ dans $E$)}
\end{aligned}
\]

Donc $h_{\mathrm{fix}}$ (et donc $\mathsf{Y}(\lambda h.E)$) se
comporte, à des $\beta$-réductions près, comme la fonction récursive
recherchée.

\medskip

\itempoint Si l'évaluation (i.e., la $\beta$-réduction) de $E$ termine
et ne fait pas intervenir $h$, alors $h_{\mathrm{fix}}$ donne juste
$E$, sinon elle itère avec $h_{\mathrm{fix}}$ pour $h$ jusqu'à ce que
ce soit le cas : c'est bien ce que fait un appel récursif.

\end{frame}
%
\begin{frame}
\frametitle{Digression / variante : le combinateur $\mathsf{Z}$}

\itempoint Le bon fonctionnement du combinateur $\mathsf{Y}$ dépend du
fait que la stratégie de $\beta$-réduction utilisée est extérieure
gauche.  Sinon le redex $(\lambda x.f(x x)) (\lambda x.f(x x))$ peut
causer un cycle de $\beta$-réductions.

\medskip

\itempoint La variante suivante évite ce problème pour définir une
fonction par appels récursifs dans un langage qui n'évalue pas « dans
les $\lambda$ » :
\[
\mathsf{Z} := \lambda f. ((\lambda x.f(\lambda v.x x v)) (\lambda x.f(\lambda v.x x v)))
\]

Cette fois :
\[
\begin{aligned}
\mathsf{Z}(\lambda h.E)
&= (\lambda f. ((\lambda x.f(\lambda v.x x v)) (\lambda x.f(\lambda v.x x v))))(\lambda h.E)\\
&\rightarrow (\lambda x.(\lambda h.E)(\lambda v.x x v)) (\lambda x.(\lambda h.E)(\lambda v.x x v))
=: h_{\mathrm{fix}}\\
&\rightarrow (\lambda h.E)(\lambda v.(\lambda x.(\lambda h.E)(x x)) (\lambda x.(\lambda h.E)(x x))v)
= (\lambda h.E)(\lambda v. h_{\mathrm{fix}} v)\\
&\rightarrow E[h\backslash \lambda v. h_{\mathrm{fix}} v]
\text{~(substitution de $\lambda v. h_{\mathrm{fix}} v$ pour $h$ dans $E$)}
\end{aligned}
\]
La forme $\lambda v. h_{\mathrm{fix}} v$ maintient $h_{\mathrm{fix}}$
non-évalué (exemple transp. suivant).

\end{frame}
%
\begin{frame}
\frametitle{Digression (suite) : exemple en Scheme}

{\footnotesize On prend ici l'exemple de Scheme pour avoir affaire à
  un langage fonctionnel non typé (le typage empêche l'implémentation
  directe du combinateur $\mathsf{Y}$ ou $\mathsf{Z}$ en OCaml ou
  Haskell).\par}

\bigskip

\itempoint Récursion sans combinateurs :

{\tt
(define proto-fibonacci\newline
\ \ (lambda (self) ; Pass me as argument!\newline
\ \ \ \ (lambda (n)\newline
\ \ \ \ \ \ (if (<= n 1) n\newline
\ \ \ \ \ \ \ \ \ \ (+ ((self self) (- n 1)) ((self self) (- n 2)))))))\newline
(define fibonacci (proto-fibonacci proto-fibonacci))\newline
(map fibonacci '(0 1 2 3 4 5 6))\newline
$\rightarrow$ (0 1 1 2 3 5 8)
\par}

\medskip

\itempoint L'idée est ici exactement celle de l'astuce de Quine : pour
m'appeler « moi-même », je m'attends à me recevoir moi-même en
argument, et je reproduis ceci lors de l'appel.

\medskip

\itempoint Le combinateur $\mathsf{Y}$ (ou $\mathsf{Z}$) automatise
cette construction.

\end{frame}
%
\begin{frame}
\frametitle{Digression (suite) : exemple en Scheme}

{\tt
;; (define y-combinator\newline
;; \ \ (lambda (f)\newline
;; \ \ \ \ ((lambda (x) (f (x x))) (lambda (x) (f (x x))))))\newline
(define z-combinator\newline
\ \ (lambda (f)\newline
\ \ \ \ ((lambda (x) (f (lambda (v) ((x x) v))))\newline
\ \ \ \ \ (lambda (x) (f (lambda (v) ((x x) v)))))))\newline
(define pre-fibonacci\newline
\ \ (lambda (fib)\newline
\ \ \ \ (lambda (n)\newline 
\ \ \ \ \ \ (if (<= n 1) n\newline
\ \ \ \ \ \ \ \ \ \ (+ (fib (- n 1)) (fib (- n 2)))))))\newline
(define fibonacci (z-combinator pre-fibonacci))\newline
\par}

\end{frame}
%
\begin{frame}
\frametitle{Représentation de l'opérateur $\mu$ de Kleene}

{\footnotesize Rappel : $\mu g(x_1,\ldots,x_k)$ est le plus petit $z$
  tel que $g(z,x_1,\ldots,x_k) = 0$ et $g(i,x_1,\ldots,x_k)\downarrow$
  pour $0\leq i<z$, s'il existe.\par}

\bigskip

\itempoint On peut faire l'algorithme « rechercher à partir
de $z$ » par appels récursifs :
\[
h(z,x_1,\ldots,x_k) = \left\{
\begin{array}{l}
z\quad\text{~si~}g(z,x_1,\ldots,x_k) = 0\\
h(z+1,x_1,\ldots,x_k)\quad\text{~(récursivement)~sinon}\\
\end{array}
\right.
\]
{\footnotesize Alors $\mu g(x_1,\ldots,x_k) = h(0,x_1,\ldots,x_k)$.\par}

\medskip

\itempoint $T := \lambda pmn.p(\lambda y.n)m$ représente $(p,m,n)
\mapsto \left\{\begin{array}{ll}m&\text{~si~}p=0\\n&\text{~si~}p\geq
1\end{array}\right.$

\itempoint $A := \lambda mfx.f(mfx)$ représente $z \mapsto z+1$

\bigskip

\itempoint La récursion est implémentée avec le
combinateur $\mathsf{Y}$ :
\[
\begin{aligned}
&\mathsf{Y}
(\lambda h z x_1\cdots x_k.
T
(v z x_1\cdots x_k)
z
(h (A z) x_1\cdots x_k)
)\\
\rightarrow\strut & \mathsf{Y}
(\lambda h z x_1\cdots x_k.
(v z x_1\cdots x_k)
(\lambda y.h (A z) x_1\cdots x_k)
z
)\\
\end{aligned}
\]

\end{frame}
%
\begin{frame}
\frametitle{Équivalence entre $\lambda$-calcul et fonctions récursives}

\itempoint Toute fonction générale récursive (i.e.,
\alert{calculable} !) $\mathbb{N}^k \dasharrow \mathbb{N}$ est
représentée par un terme du $\lambda$-calcul (sous les conventions
données : application aux entiers de Church, réduction extérieure
gauche).

\bigskip

\itempoint Réciproquement, toute fonction $\mathbb{N}^k \dasharrow
\mathbb{N}$ représentable par un terme du $\lambda$-calcul est
calculable, car on peut implémenter la réduction extérieure gauche.

\bigskip

\itempoint Bref, $f\colon \mathbb{N}^k \dasharrow \mathbb{N}$ est
représentable par un terme du $\lambda$-calcul \alert{ssi} elle est
générale récursive.

\bigskip

\itempoint De plus, cette équivalence est \alert{constructive} : il
existe des fonctions p.r. :
\begin{itemize}
\item l'une prend en entrée le numéro $e$ d'une fonction générale
  récursive (et l'arité $k$) et renvoie le code d'un terme du
  $\lambda$-calcul qui représente cette $\varphi_e^{(k)}$,
\item l'autre prend en entrée le code d'un terme du $\lambda$-calcul
  qui représente une fonction $f$, et son arité $k$, et renvoie un
  numéro $e$ de $f$ dans les fonctions générales récursives $f =
  \varphi_e^{(k)}$.
\end{itemize}

\end{frame}
%
\section{Conclusion}
\begin{frame}
\frametitle{Récapitulation}

\itempoint\textbf{Théorème} : les fonctions $\mathbb{N}^k \dasharrow
\mathbb{N}$ \textbf{(1)} générales récursives,
\textbf{(2)} représentables en $\lambda$-calcul, et
\textbf{(3)} calculables par machine de Turing, coïncident toutes.

\bigskip

\centerline{On les appelle les \textbf{fonctions calculables}.}

\bigskip

\itempoint De plus, ces équivalences sont constructives : on peut
passer algorithmiquement (= calculablement !) d'une représentation à
l'autre.

{\footnotesize Ce sont des formes de \alert{compilation} d'un langage
  en un autre.\par}

\bigskip

\itempoint Les questions suivantes \alert{ne sont pas décidables}
algorithmiquement :
\begin{itemize}
\item savoir si une machine de Turing donnée s'arrête,
\item savoir si un terme du $\lambda$-calcul est normalisable.
\end{itemize}

\end{frame}
%
\begin{frame}
\frametitle{Turing-complétude}

Un langage de programmation est dit \textbf{Turing-complet} lorsque
(convenablement idéalisé !) il permet d'implémenter précisément les
fonctions calculables au sens de Church-Turing.

\bigskip

Un ordinateur réel ne peut \alert{jamais faire plus} qu'une machine de
Turing (sauf p.-ê. : faire du hasard vrai).  La question est de
savoir si le langage permet \alert{autant}.

\bigskip

Tous les langages de programmation généralistes \alert{sont
  Turing-complets} : Python, Java, JavaScript, C, C++, OCaml, Haskell,
Lisp, Perl, Ruby, Smalltalk, Prolog…

\bigskip

Certains le sont même plus ou moins « par accident » : CSS, TeX, XSLT, m4…
(parfois sous conditions, ou sous réserve d'interprétation).

\bigskip

Pas toujours clair : assembleurs (pas évident d'idéaliser les entiers).

\bigskip

Conséquence du problème de l'arrêt : on ne peut pas algorithmiquement
décider si un programme donné (en Python, etc.) termine ou non.

\end{frame}
%
\begin{frame}
\frametitle{Turing tarpit}

{\footnotesize « Fosse à bitume de Turing » ?\par}

\bigskip

\itempoint Tous les langages usuels se valent du point de vue de la
calculabilité.  Ce n'est pas pour autant qu'ils se valent en
pratique !  (En commodité et/ou efficacité.)

\bigskip

\itempoint Ça ne signifie pas qu'un langage Turing-complet peut
forcément « tout » faire.  Par exemple, un langage qui ne permet comme
entrée/sortie que d'afficher des entiers peut être Turing-complet et
ne permet pas d'écrire « bonjour ».

\bigskip

\itempoint Si en principe on peut convertir toute fonction calculable
dans tout langage Turing-complet, la conversion peut devenir
extrêmement inefficace, malcommode ou illisible.

\end{frame}
%
\begin{frame}
\frametitle{$\lambda$-calcul non typé et type récursif}

{\footnotesize\textcolor{gray}{Remarque faite pour plus tard.}\par}

\smallskip

Le fait d'avoir un type $t$ tel que $t \,\cong\, (t\rightarrow t)$
permet d'implémenter dans ce type le $\lambda$-calcul « non typé »
(donc tue l'espoir de décider la terminaison).

\bigskip

Exemple en OCaml (ici, \texttt{loop} produit une boucle infinie) :

\smallskip

{\footnotesize
\texttt{type t = T of (t -> t)}\newline
\texttt{let apply : t -> t -> t = fun (T rator) -> fun rand -> rator rand}\newline
\texttt{let id : t = T (fun x -> x)}\hfill\texttt{(* }$\lambda x.x$\texttt{ *)}\newline
\texttt{let ch0 : t = T (fun f -> T (fun x -> x))}\hfill\texttt{(* }$\lambda fx.x$\texttt{ *)}\newline
\texttt{let ch1 : t = T (fun f -> T (fun x -> apply f x))}\hfill\texttt{(* }$\lambda fx.fx$\texttt{ *)}\newline
\texttt{let ch2 : t = T (fun f -> T (fun x -> apply f (apply f x)))}\hfill\texttt{(* }$\lambda fx.f(fx)$\texttt{ *)}\newline
\texttt{let om : t = T (fun x -> apply x x)}\hfill\texttt{(* }$\lambda x.xx$\texttt{ *)}\newline
\texttt{let loop : t = apply om om}\hfill\texttt{(* }$(\lambda x.xx)(\lambda x.xx)$\texttt{ *)}\newline
\texttt{(* let loop = (fun (T h) -> h (T h)) (T (fun (T h) -> h (T h))) *)}\newline
\par}

\medskip

Remarquer qu'ici on arrive à provoquer une boucle infinie sans aucun
\texttt{let rec} (et malgré le typage).

\end{frame}
%
\begin{frame}
\frametitle{Une méditation googologique}

{\footnotesize\textcolor{gray}{Ceci est une sorte de digression, pour
    inviter à la réflexion.}\par}

\medskip

{\footnotesize « googologie » = étude des grands nombres ; de
  « googol », nom fantaisiste de $10^{100}$\par}

\bigskip

On cherche à minorer calculabl\textsuperscript{t} la fonction « castor
affairé », c-à-d :
\begin{itemize}
\item concevoir un programme dans un langage de programmation idéalisé
  (machine de Turing, $\lambda$-calcul, Python, OCaml…),
\item de taille « humainement raisonnable » (peu importent les détails),
\item qui \alert{termine en temps fini} (théoriquement !),
\item mais calcule un nombre aussi grand que possible (variante :
  attend un temps aussi long que possible).
\end{itemize}

\bigskip

Exemple : implémenter $A_\Delta\colon n \mapsto A(n,n,n)$ (fonction
d'Ackermann diagonale) et calculer $A_\Delta(A_\Delta(\cdots(100))) =
A_\Delta^{\circ 100}(100)$ ou qqch du genre.

…On peut faire \textcolor{orange}{beaucoup} plus grand !

\end{frame}
%
%% \mode<article>{
%%     \ifodd\value{page}
%%         \clearpage\hbox{}\newpage
%%         \thispagestyle{empty}
%%     \fi}
%
\end{document}
